\chapter{Narendra Avasthi: Complete Graded Problem Bank (Levels 1, 2 \ and  3)}
\label{chap:narendra_avasthi_all_questions}

\begin{problembox}
\textbf{Problem 1} \hfill \textbf{[Narendra Avasthi Q1]}
\label{q:ext-narendra_avasthi_all_questions-1}

NaCl(Molten)withPtelectrode
Cl2(g)
Na(l)
\end{problembox}

\begin{problembox}
\textbf{Problem 2} \hfill \textbf{[Narendra Avasthi Q2]}
\label{q:ext-narendra_avasthi_all_questions-2}

NaCl(aq)withPtelectrode
Cl2(g)
H(g)
\end{problembox}

\begin{problembox}
\textbf{Problem 3} \hfill \textbf{[Narendra Avasthi Q3]}
\label{q:ext-narendra_avasthi_all_questions-3}

Na2SO4(aq)with Pt electrode
02(g)
H2(g)
\end{problembox}

\begin{problembox}
\textbf{Problem 4} \hfill \textbf{[Narendra Avasthi Q4]}
\label{q:ext-narendra_avasthi_all_questions-4}

NaNO3(aq)withPtelectrode
O2(g)
H2(g)
\end{problembox}

\begin{problembox}
\textbf{Problem 5} \hfill \textbf{[Narendra Avasthi Q5]}
\label{q:ext-narendra_avasthi_all_questions-5}

AgNO3(aq)withPtelectrode
O2(g)
Ag(s)
\end{problembox}

\begin{problembox}
\textbf{Problem 6} \hfill \textbf{[Narendra Avasthi Q6]}
\label{q:ext-narendra_avasthi_all_questions-6}

[CuSOq(aq)with inert electrode
O2(g)
Cu(s)
\end{problembox}

\begin{problembox}
\textbf{Problem 7} \hfill \textbf{[Narendra Avasthi Q7]}
\label{q:ext-narendra_avasthi_all_questions-7}

CuSO4(aq)with copper electrode
Cu dissolve
Cu(s)
Galvanic Cell
This cell converts chemical energy into electrical energy.
Anodeof Zn
Salt
Cathode of Cu
Bridge
(-ve electrode)
(+ve electrode)
Cotton
(Anodiccelil)
plug
(Cathodic cell)
Galvanic cell is made up of two half cells i.e., anodic and cathodic. Oxidation takes place at
anode and reduction at cathode. It is also known as voltaic cell. It may be represented as shown in
Fig.Zinc rod immersed in ZnSO4behaves as anode and copper rod immersed inCusObehaves as
cathode.
Anode half cell reaction:
Zn(s)---$/to$Zn$2$+(aq)+ 2e(oxidation)
Cathodehalfcellreaction:
Cu$2$+(aq)+2e---$/to$Cu(s)(reduction)
Overallprocess:
Zn(s)+Cu$2$(aq)---$/to$$/to$Cu(s)+Zn$2$+(aq)
In galvanic cel like Daniel cel; electrons flow from anode (zinc rod) to the cathode (copper
rod) through extermal circuit; zinc dissolves as Zn2+;Cu2+ ion in the cathode cell picks up two
electron andbecomedeposited atcathode.
Salt Bridge
It is U-shaped tube contains saturated solution of inert electrolyte likeKCI,KNO3,NH4Cl and
NH4NO, etc. in agar-agar gel.
Cation and anion of inert electrolyte have same mobility.
lons of inert electrolyte do notmix withelectrolytic solution present in the half cell.
Ions of inert electrolyte do not participate in electrochemical change.
--- PAGE 3 ---
390
PROBLEMSIN CHEMISTRY
Function of Salt Bridge
Itcomplete theelectricalcircuit.
It maintained thetwo half cell electricallyneutral by theflow of ions.
Representation of a cell (TUPAC conventions): Let us illustrate the convention taking
the example of Daniel cell.
.:.Zn(s)|ZnSO4(aq)l|CuSO4(aq)|Cu(s)
(i)Twohalf cells are separated by double vertical lines and it indicate salt bridge or any type
of porous partition.
(ii) EMF (electromotive force) may be written on the right hand side of the cell.
(iv) Single vertical lines indicate the phase separation between electrode and electrolyte
solution.
(s)n(b)+nl(bo)+uz ()uz
(-)
(+)
Anode
Catode
Eeel = EOP(LHs) +ERr(RHs) Ecl =EoP(LH) +ERP(RHs) Ecel =ERP(RHs) -ERP(LHS)
Nernst's Equation
Ecell = Ecel
2.303RT
.logQ
nF
At 298 K temp.
Ecell =E
0.0591
logQ ;where Q is reaction quotient.
cell
n
Cell Thermodynamics
(i) $/Delta$G = -nFEcel
aut =09V (1).
(i) Equilibrium constant (K), log K=
0.0591
(iv) Temp. coefficient of cell: 
OE
aT)P
(v)Enthalpy of reaction inside the cell, $/Delta$H = nFE + nFT
OE
aT
Concentration Cell
Anode and cathode both are formed by same substance so Ece = O, such type of celis known as
concentration cell.Concentration gradient may arise either in electrodematerial or ini electrolyte.
Thus there are two types of concentration cell.
Electrode gas conc. cell :
Pt,H(g)|H+(C)|H(g),Pt
B atm
P2 atm
0.0591log
Ecell
P
log
2
P2
--- PAGE 4 ---
ELECTROCHEMISTRY
Electrolyte conc. cell :
Zn(s)|ZnSO4(C)I|ZnSO4(C2)Zn(s)
0.0591
C2
log
2
C
DifferentType,ofElectrode
S.No.
Name of Electrode
Anode
Cathode
\end{problembox}

\begin{problembox}
\textbf{Problem 8} \hfill \textbf{[Narendra Avasthi Q1]}
\label{q:ext-narendra_avasthi_all_questions-8}

Hydrogen electrode
Pt(s)|H2(g)|H+(aq)
H+(aq)|H2(g)/Pt(s)
0.0591
[H+)2
0.0591
PH2
1og
2
PH2
2
[H+]2
\end{problembox}

\begin{problembox}
\textbf{Problem 9} \hfill \textbf{[Narendra Avasthi Q2]}
\label{q:ext-narendra_avasthi_all_questions-9}

Metal-metalionelectrode
M (s)|M+n(aq)
M+n(aq)| M (s)
0.0591
log[M+n]
0.0591
1
log
n
n
[M+n]
\end{problembox}

\begin{problembox}
\textbf{Problem 10} \hfill \textbf{[Narendra Avasthi Q3]}
\label{q:ext-narendra_avasthi_all_questions-10}

Calomelelectrode
Hg (l), Hg2Cl2(s)/ CI-(aq)
Cl (aq) | Hg2Cl2(s), Hg (l)
0.0591
1
0.0591
log
log[C1-]$2$
2
[C1-]$2$
2
\end{problembox}

\begin{problembox}
\textbf{Problem 11} \hfill \textbf{[Narendra Avasthi Q4]}
\label{q:ext-narendra_avasthi_all_questions-11}

Redoxelectrode
Pt(s)/Fe+2(aq), Fe+3(aq)
Fe+3(aq), Fe+2(aq)|Fe(s)
0.0591
[Fe+3]
ERP
0.0591.
[Fe+2]
log
1
[Fe+2]
1
[Fe+2]
\end{problembox}

\begin{problembox}
\textbf{Problem 12} \hfill \textbf{[Narendra Avasthi Q5]}
\label{q:ext-narendra_avasthi_all_questions-12}

Metal insoluble salt anion
Ag (s)| AgCl(s)| Cl-(aq)
Cl"(aq)| AgCl(s)| Ag (s)
electrode
0.0591
1
0.0591
log
log[Cl-]
1
[CI-]
1
Batteries
Electrochemical cells can be used as batteries.Batteries are of two kinds:
Primary battery:Where the reaction occurs only once and can not be reused once it
becomes dead over the course of time.For examples, dry cell (Leclanche cell),Mercury cell.
Secondary battery:Which can be recharged by passing current throughit in the opposite
direction so that it canbeused again.For example,lead-acid battery,Nickel-cadmium cell.
Fuel Cells
methane,methanol,etc.)directly into electrical energy.
Fuel cells aremore eficient thermodynamicallyandmore of theenergyof thereactioncanbe
made available for useful work provided that the supply of reactants is maintained.
Methanol,ethanol,hydrazine,formaldehyde,carbonmonoxide canbe used asfuels infuelcells
apartfrom hydrogen.
--- PAGE 5 ---
0
PROBLEMSINCHEMISTRY
Corrosion
When the metal exposed to some environment it gets converted to its oxides. The oxidative
deterioration of metalisknown as corrosion.
Example:Rusting of iron, tarnishing of silver, development of green coating on copper and
bronze etc.
Oxidation
Fe(s)---/toFe2+(aq)+2e
Reduction
O2(g)+ 4H+(aq) +4e/textasciitilde()-/to 2H2O(l).
Atmospheric oxygen oxidised Fe2+
2Fe2+(aq)+ 2H2O(l)+O2(g)---/toFe2O3(s)+4H*(aq)
2
Hydrated feric oxide (Fe2O3/cdotxH2O(s)) is known as rust.
Factors which enhance corrosion :
i)
Presence of impurities in the metal.
(i) Presence of moisture.
(i) Presence of electrolyte.
Prevention of corrosion:
(i)'Barrier protection by oil/grease layer, paints or electroplating.
(i)Sacrificial protection by coating the metal with more electropositive metal.
Conductance
Conductance (G):It is defined as the reciprocal of the electrical resistance i.e.,G =1/R.It
measures the easewith which the current flows through a conductor.
Unit : Siemen, S or Q-1.
Specific resistance or resistivity (o):Theresistivity or specific resistance is defined as
theresistance in /Omega of a conductor havinglength equalto1 cm andarea of cross-section
equal to 1 cm2.
orR=p
R
wherep= specificresistance,l=length of conductor
A
A
A =area of cross-section of the conductor.
Unit:/Omega cm.
W
Specific conductance or conductivity (k):It is define as the reciprocal of specific
resistance.
i.e.,
k=
1
:R=p
1
1
k=Gx
p
A
A
k=GxG*
where G*=cell constant
A
When I =1 cm and A =1 cm2; k =G
Thus conductivity is the conductance of one centimeter cube or conductance of one cm
cube of the solution of an electrolyte.
Unit : Q-'cm -', Scm -1.
Molar conductivity (Am):The conducting power of all the ions produced by dissolving 1.
mole of an electrolyte in solution.
--- PAGE 6 ---
393
ELECTROCHEMISTRY
Mathematically, /LambdaM=kxV
AM(Scm2 mol-1) =k/times 1000
(M = molarity)
M
Variation of conductivity and molar conductivity with concentration:
Forweak andstrongelectrolytes
Conductivity decreases with concentration.
Molarconductivityincreaseswithdecreaseinconcentration.
01
/textasciicircum()m
Strongelectrolyte
(NaCI)
Weakelectrolyte
/C
Kohirausch's Law.
According to this law,at infinite dilution,when' the dissociation is complete,each ion makes a
definite contribution towards molar conductivity of theelectrolyteirrespective of the nature of the
otherions present.
Theimolarconductivity of anelectrolyte at infinite dilutionis the sum of theionic conductivities
of the cations and the anions each multiplied by the number of ions present in one formula unit of
the electrolyte e.g.,ABy.
The equivalent conductivity of an electrolyte.atinfinite dilutionis the sum of the equivalent
conductivities of the cations and anions.
(AxB)=xx(A+)+y(Bx-)
Nq.(AxB)=2g.(A+)+2g.(Bx-)
Application
(a) Determination of equivalent/molar conductivities of weak electrolytes at
infinite dilution
A/%(CHCOOH)=A(CHCOONa)+A/%(HCI)-A/%(NaCI)-A/%(NaCI)
(b)Determination of degree of dissociation(a) and equilibrium constant(K)of
an electrolyte ata given dilution
MolarconductanceatconcentrationC
M
/alpha=
Molarconductanceatinfinitedilution
WV
HAH+A
Initial conc.
c.
Conc.at equilibrium C-Ca
Ca
Ca
--- PAGE 7 ---
PROBLEMSINCHEMISTRY
Ca
/LambdaM
(c)Determinationof thesolubility ofasparingly solublesalt
Since the solution is saturated at infinite dilution A = A and molarity = solubility.
/cdotA/%=k/timesx1000
molarity
Solubility(S)=
kx1000
For sparingly soluble salt AxB, : Solubility product Ksp =x* y . sx+y
ConductometricTitration
the titresolution(forburette)iscalledconductometrictitration:
Titration of HCl(aq) with NaOH(aq) (Strong base.Vs Strong acid)
HCl(aq)+NaOH(aq)----/toNaCl(aq)+H2O(l)
Endpoint
Volume.ofNaOH.added
Titration of CH;COOH(aq) with NaOH(aq) (Weak acid Vs Strong base)
CHCOOH(aq)+NaOH(aq)-/toCHCOONa+H2O.(l)
Conductance
Endpoint
VolumeofNaOH.added
--- PAGE 8 ---
ELECTROCHEMISTRY
395
Level
1.A cell reaction would be spontaneous if the cell potential and /Delta,G arerespectively:
(a) positive and negative
(b)negative,negative
(c) zero,zero
(d) positive,zero
2.Which of thefollowing statement is correct?
(a) Cathode is-ve terminal in both,galvanic and electrolytic cells
(b)Anode is+ve terminal in both,galvanic and electrolytic cells
(c) Cathode and anode are -ve terminal in electrolytic and galvanic cell.
(d) Cathode and anode are +ve terminal in electrolytic and galvanic cell.
3.Electrolytes when dissolved in water dissociate into ions because:
(a)They are unstable.
(b) The water dissolves it.
(c)The force ofrepulsion increases.
(d)Theforce ofelectrostatic atractionisbroken down bywater.
4.The electric charge required for electrode deposition of one gram-equivalent of a substnace is:
(a) one ampere per second
(b) 96500 coulombs per second
(c)one ampere for onehour
(d) charge on one mole of electrons
\end{problembox}

\begin{problembox}
\textbf{Problem 13} \hfill \textbf{[Narendra Avasthi Q5]}
\label{q:ext-narendra_avasthi_all_questions-13}

The amount of an ion liberated on an electrode during electrolysis does not depend upon:
(a) Conductance of the solution
(b) Current strength
(c)Time
(d)Electrochemical equivalentof theelement
\end{problembox}

\begin{problembox}
\textbf{Problem 14} \hfill \textbf{[Narendra Avasthi Q6]}
\label{q:ext-narendra_avasthi_all_questions-14}

How many electrons are there in one coulomb of electricity?
(a) 6.023$/times$1023
(b) 1.64$/times$ 10-24
(c) 6.24 x 1018
(d) 6.24 $/times$ 10-24
7.How many coulombs are provided by a current of 0.010 mA in the calculator battery that can
operate for 1000 hours?
(a) 1.0
(b)10
(c)0.010
(d)36
8.How many minutes are required to deliver 3.21 $/times$ 10$/circ$ coulombs using a current of 500 A used
in the commercial productionof chlorine?
(a) 8.3
: (b) 5.3x104
(c)6420
(d) 107
9.Passage of a current for 548 seconds through a silver coulometer results in the deposition of
0.746 g of silver. What is the current (in A)?
(a) 1.22
(b) 1.16
(c) 1.07
(d) 1.00
10.Electrolysis can be used to determine atomic masses.A current of 0.550 A deposits 0.55 g ofa
certain metal in100 minutes.Calculate the atomic mass of the metal if eq.wt.=mole.wt./3
(a) 100
(b) 45.0
(c)48.25
(d) 144.75
11.Beryllium occurs naturally in the form of beryl. The metal is produced from its ore by
electrolysis after the ore hasbeen converted to the oxide and then to the chloride.Howmany
grams,of Be(s) is deposited from a BeCl2 solution by a current of 5.0 A that flows for 1.0 h?
(Atomic weight:Be=9)
(a)0.840
(b) 1.68
(c)1.42
(d) 1.08
--- PAGE 9 ---
396
PROBLEMSINCHEMISTRY
12.How many minutes will it take to plate out 5.0g of Cr from a Cr2(SO4) solution using a
current of 1.50 A? (Atomic weight : Cr = 52.0)
(a) 254
(b)30g
(c)152
(d)103
13.Calculate the current (in mA) required to deposit 0.195 g of platinum metal in 5.0 hours from
a solution of PtCl2 :(Atomic weight: Pt = 195)
(a)310
(b) 31
(c) 21.44
(d) 5.36
\end{problembox}

\begin{problembox}
\textbf{Problem 15} \hfill \textbf{[Narendra Avasthi Q14]}
\label{q:ext-narendra_avasthi_all_questions-15}

How many Faradays are required to reduce 0.25 g of Nb (V) to the metal?
(Atomic weight : Nb = 93g)
(a) 2.7$/times$10-3
(b) 1.3 $/times$ 10-2
(c) 2.7 $/times$10-2
(d) 7.8 x 10-3
\end{problembox}

\begin{problembox}
\textbf{Problem 16} \hfill \textbf{[Narendra Avasthi Q15]}
\label{q:ext-narendra_avasthi_all_questions-16}

One gm metal M3+ was discharged by the passage of 1.81 $/times$ 1023 electrons. What is the atomic
weight of metal?
(a) 33.35
(b) 133.4
(c) 66.7
(d)None of these
16.Total charge required for the oxidation of two moles MnsO4 into MnO2 in presence of
aikaline medium is :
(a)5F
(b) 10 F
(c)20 F
(d)None of these
17.The electrolytic decomposition of dilute sulphuric acid with platinum electrode, cathodic
reaction is :
(a) Reduction of H+(b) Oxidationof SO$2$(c) Reduction SO$2$(d) Oxidation of H2O
18.Which one of the following metals can not be obtained on electrolysis of aqueous solution of
its salts ?
(a) Mg
(b) Ag
(c) Cu
(d) Cr
19.A solution of potassium sulphate in water is electrolysed using inert electrodes.The products
at the cathode and anode are respectively.
(a)H2,O2
(b) O2,H2
(c) O2, Na
(d) None of these
20.The passage of current through a solution of certain$/cdot$electrolyte results in the evolution of
H2(g) at cathode and Cl2(g) at anode. The electrolytic solution is :
(a)Water
(b) aq.H2SO4
(c) aq. NaCl
(d) aq. CuCl2
\end{problembox}

\begin{problembox}
\textbf{Problem 17} \hfill \textbf{[Narendra Avasthi Q21]}
\label{q:ext-narendra_avasthi_all_questions-17}

When an aqueous solution of H2SO4 is electrolysed, the product at anode is :
(a) H"
(b) OH-
(c) SO$2$
(d) O2
\end{problembox}

\begin{problembox}
\textbf{Problem 18} \hfill \textbf{[Narendra Avasthi Q22]}
\label{q:ext-narendra_avasthi_all_questions-18}

An aqueous solution of Na2SO4 in water is electrolysed using Pt electrodes.The products at
the cathode and anode are respectively:
(a)H2,SO2
(b)O2,NaOH
(c)H2,O2
(d) O2,SO2
23.The electrolysis of a solution resulted in the formation of H2(g)at the cathode and O2(g)at
the anode.The solution is :
(a)AgCl(aq)
(b) H2SO4(aq)
(c)highly concentrated NaCl(aq) solution
(d) CuCl2(aq)
24.If mercury is used as cathode in the electrolysis of aqueous NaClsolution,the ion discharged at
cathode is :
(a)H+
(b) Na*
(c) OH"
(d) C1-
\end{problembox}

\begin{problembox}
\textbf{Problem 19} \hfill \textbf{[Narendra Avasthi Q25]}
\label{q:ext-narendra_avasthi_all_questions-19}

A dilute aqueous solution of CusO4 is electrolyzed using platinum electrodes.The products at
the anode and cathode are :
(a) O2, H2
(b) H2,O2
(c) O2, Cu
(d) S2O , H2
--- PAGE 10 ---
397
ELECTROCHEMISTR
\end{problembox}

\begin{problembox}
\textbf{Problem 20} \hfill \textbf{[Narendra Avasthi Q26]}
\label{q:ext-narendra_avasthi_all_questions-20}

What products are formed during the electolysis of concentrated aqueous solution of sodium
chloride?
(1) Cl2(g) at anode
(I) NaOH as electrolyte
(II) H2(g) at cathode
(a) I only
(b)I and II only
(c)I and l only
(d)1, II and III
\end{problembox}

\begin{problembox}
\textbf{Problem 21} \hfill \textbf{[Narendra Avasthi Q27]}
\label{q:ext-narendra_avasthi_all_questions-21}

Which of the following aqueous solution produces metal after electrolysis ?
(a)K2Cr2O7
(b) KMnO4
(c)CHCOONa
(d) CuCl2
\end{problembox}

\begin{problembox}
\textbf{Problem 22} \hfill \textbf{[Narendra Avasthi Q28]}
\label{q:ext-narendra_avasthi_all_questions-22}

How much time is required for complete decomposition of 4 moles of water using 4 ampere?
(a) 3.86$/times$ 10 sec(b) 1.93$/times$105 sec(c) 96500 sec
(d) 48250 sec
\end{problembox}

\begin{problembox}
\textbf{Problem 23} \hfill \textbf{[Narendra Avasthi Q29]}
\label{q:ext-narendra_avasthi_all_questions-23}

An aqueous solution containing 1 M each of Au3+,Cu2,Ag*, Li is being electrolysed by
using inert electrodes. The value of standard potentials are:
E$/circ$
Ag/Ag
=0.80V,E$/circ$
Cu+/Cu
With increasing voitage, the sequence of deposition of metals on the cathode will be:
(a)Li, Cu, Ag,Au(b) Cu,Ag,Au
(c)Au, Ag, Cu
(d) Au, Ag, Cu, Li
30.If 0.50 L of a 0.60 MSnSO4 solution is electrolyzed for a period of 30.0 min using a current of
4.60 A. If inert electrodes are used, what is the final concentration of Sn? remaining in the
solution?[at.wt.of Sn =119]
(a)0.342M
(b)0.544 M
(c)0.389 M.
(d) 0.514 M
31.A 100.0 mL dilute solution of Ag+ is electrolyzed for 15.0 minutes with a current of 1.25 mA
and the silver is removed completely.What was theinitial [Ag+]?
(a) 2.32$/times$ 10-1
(b) 2.32$/times$104
(c)2.32$/times$10-3
(d) 1.16x 10-4
\end{problembox}

\begin{problembox}
\textbf{Problem 24} \hfill \textbf{[Narendra Avasthi Q32]}
\label{q:ext-narendra_avasthi_all_questions-24}

A 250.0 mL sample of a 0.20 M Cr3+ is electrolyzed with a current of 96.5 A. If the remaining.
[Cr*+] is 0.1 M the duration of process is :
(a) 25 sec
(b) 225 sec
(c) 150 sec
(d) 75 sec
33.The element indium is tobe obtained by electrolysis of amolten halide of the element.Passage
of a current of 3.20 A for a period of 40.0 min results in formation of 3.05 g of In.What is the
oxidation state of indium in the halide melt?(Atomic weight:In=114.8)
(a) 3
(b) 2
(c) 5
(d) 1
\end{problembox}

\begin{problembox}
\textbf{Problem 25} \hfill \textbf{[Narendra Avasthi Q34]}
\label{q:ext-narendra_avasthi_all_questions-25}

An electrolysis of a oxytungsten complex ion using 1.10 A for 40 min produces 0.838 g of
tungsten.What is the charge on tungsten in the material?(Atomic weight:W =184)
(a) 6
(b) 2
(c) 4
(d) 1
35.In the electrolysis of aqueous NaCl,what volume ofCl2(g)is produced in the time that it takes
to liberate 5.o liter of H2(g)?Assume that both gases are measured at STP
(a) 5.0
(b) 2.50
(c) 7.50
(d) 10.0
36.How many grams of Cr are deposited in the electrolysis of solution of Cr(NO3)3 in the same
time that it takes to deposit 0.54g of Ag in a silver coulometer arranged in series with the
.Cr(NO3) cell?(Atomic weight:Cr=52.0;Ag=108)
(a) 0.0866
(b)0.0288
(c)0.173
(d)0.220
37.In the electrolysis of aCuSO4 solution,how many grams of Cu are plated out on the cathode in
the time that it takes to liberate 5.6 litre of O2(g), measured at STP at the anode?
(a) 31.75
(b) 14.2
(c) 4.32
(d)None of these
38.Ammonium perchlorate,NH4ClO4,used in the solid fuel in the booster rockets on the space
shuttle,isprepared from sodium perchlorate,NaClO4,whichis produced commerciallybythe
--- PAGE 11 ---
398
PROBLEMSINCHEMISTRY
electrolysis of a hot,stirred solution of sodium chloride.Howmanyfaradays are required to
produce 1.0kg of sodium perchlorate?
NaCl+4H2O-$/to$NaClO+4H2
(a)40.3
(b) 18.3
(c) 31.6
(d) 65.3
\end{problembox}

\begin{problembox}
\textbf{Problem 26} \hfill \textbf{[Narendra Avasthi Q39]}
\label{q:ext-narendra_avasthi_all_questions-26}

In the commercial preparation of aluminum, aluminum oxide (Al2O3) is electrolyzed at
1000$/circ$C.How many coulombs of electricity arerequired togive54kgof aluminum?Assume
following reaction takes place at cathode:
++
(a) 17.3x 108
(b) 3.21x 107
(c) 1.82 $/times$104
(d) 57.9x 107
40.When molten lithium chloride (LiCl) is electrolyzed, lithium metal is formed at the cathode.If
current efficiency is 75/% then how many grams oflithium are liberated when1930C of
charge pass through the cell? (Atomic weight : Li = 7)
(a) 0.105
(b) 0.120
(c) 0.28
(d) 0.240
41.Sodium metal is produced commercially by the electrolysis of molten sodium chloride and
chlorine is produced as a by product.How many litres of chlorine at 1.8 atm and 27$/circ$C will be
produced if a current of 1.0x 103 A is passed through NaCl (l) for 9.65 h?
(a)2463
(b)460
(c)1800
(d) 1231.6
42.H2(g)andO2(g),can be produced by the electrolysis of water.What total volume (in L) ofO2
and H2 are produced at STP when a current of 30 A is passed through a K2SO4 (aq) solution
for 193 min.?
(a) 20.16.
(b) 40.32
(c) 60.48
(d) 80.64
\end{problembox}

\begin{problembox}
\textbf{Problem 27} \hfill \textbf{[Narendra Avasthi Q43]}
\label{q:ext-narendra_avasthi_all_questions-27}

The cost of 2 Rs/kWh of operating an electric motor for 10 hours takes 10 amp at 110V is:
(a) 79200 Rs
(b)22000 Rs
(c)220 Rs
(d) 22 Rs
\end{problembox}

\begin{problembox}
\textbf{Problem 28} \hfill \textbf{[Narendra Avasthi Q44]}
\label{q:ext-narendra_avasthi_all_questions-28}

A 1 M solution of H2SO4 is electrolyzed. Select right statement with products at anode and
cathoderespectively:
Given:
2SO$2$---$/to$S202+ 2e";E$/circ$=-2.01 V
H2O(l)-$/to$2H+(aq)+1/2O2(g)+2e;E$/circ$=-1.23V
(a) concentration of H2SO4 remain constant; H2,O2
(b) concentration of HSO4increases;O2,H2
(c) concentration of H2SO4 decreases;O2,H2
(d) concentration of H2SO4remains constant; S2O,H2
45.Cadmium amalgam is prepared by electrolysis of a solution of CdClusing a mercury cathode.
How long should a current of 4 A be passed in order to prepare 10/% by wt. Cd in Cd---Hg
amalgam on cathode of 4.5 g Hg? (atomic wt. of Cd = 112)
(a) 400 sec
(b) 215.40 sec
(c)861.6 sec
(d) 430.8 sec
\end{problembox}

\begin{problembox}
\textbf{Problem 29} \hfill \textbf{[Narendra Avasthi Q46]}
\label{q:ext-narendra_avasthi_all_questions-29}

Use of electrolysis is in :
(a)Electrorefining
(b)Electroplating
(c) Both(a)and (b)
(d)None of these
\end{problembox}

\begin{problembox}
\textbf{Problem 30} \hfill \textbf{[Narendra Avasthi Q47]}
\label{q:ext-narendra_avasthi_all_questions-30}

When a solution of AgNO3 (1 M) is electrolyzed using platinum anode and copper cathode.
What are the products obtained at two electrodes?
=+0.0volt;
H+|H2
--- PAGE 12 ---
ELECTROCHEMISTRY
399
E$/circ$
(a) Cu -$/to$ Cu$2$+ at anode; Ag+$/to$Ag at cathode
(b) H20 ---$/to$O2 at anode; Cu2+
$/to$Cuatcathode
(c)H2O---$/to$O2at anode;Ag+-$/to$Ag at cathode
(d) NO3---$/to$NO2 at anode;Ag+---$/to$Ag at cathode
\end{problembox}

\begin{problembox}
\textbf{Problem 31} \hfill \textbf{[Narendra Avasthi Q48]}
\label{q:ext-narendra_avasthi_all_questions-31}

Which of the following statements is correct about Galvanic cell ?
(a)It converts chemical energy into electrical energy.
(b)Itconvertselectricalenergy into chemicalenergy.
(c) It converts metal from its free state to the combined state.
(d)It converts electrolyte into individual ions.
49.E$/circ$forCl2(g)+2e----$/to$2Ci-(aq)is 1.36 V;E$/circ$forCl-(g)---$/to$1/2Cl2(g)+e.is :
(a):1.36 V
(b) -1.36 V
(c)-0.68V
(d) 0.68 V
\end{problembox}

\begin{problembox}
\textbf{Problem 32} \hfill \textbf{[Narendra Avasthi Q50]}
\label{q:ext-narendra_avasthi_all_questions-32}

When two half-cells of electrode potential of E, and E2 are combined to form a cell of
electrode potential E3, then
(when n1,n2 and ng are no.of electrons exchanged in first,second and combined half-cells) :
En+E2n2
(a)E3=E2-E1
(b) E3=
n3
En-E2n2
(c)E3=
(d) E=E1+E2
n
51.The function of a salt bridge is to :
(a)maintain electrical neutrality of both half cels
(b)increase the cell potential at the positive electrode
(c)decrease the cell potential at the negative electrode
(d) eliminate the impurities present in the electrolyte
\end{problembox}

\begin{problembox}
\textbf{Problem 33} \hfill \textbf{[Narendra Avasthi Q52]}
\label{q:ext-narendra_avasthi_all_questions-33}

Saturated solution of KNO with agar-agar is used to make'salt bridge'because :
(a)size of K+is greater than that of NO3
(b) velocity of NO is greater than that of K+
(c)velocities of K+ and NO are nearly the same
(d)both velocity and sizes of K+ and NO ions are same
53.A salt bridge contains :
(a) A saturated solution of KCl and agar-agar
(b)A saturated solution of KNO3 and agar-agar
(c) A saturated solution of NH4NO and agar-agar
(d) All of these
\end{problembox}

\begin{problembox}
\textbf{Problem 34} \hfill \textbf{[Narendra Avasthi Q54]}
\label{q:ext-narendra_avasthi_all_questions-34}

The nature of curve of Ecen vs. log K.  is :
(a)straightline
(b) parabola
(c)hyperbola
(d)elliptical curve
\end{problembox}

\begin{problembox}
\textbf{Problem 35} \hfill \textbf{[Narendra Avasthi Q55]}
\label{q:ext-narendra_avasthi_all_questions-35}

Consider the following equations for a cell reaction
A+B$/to$C+D;E$/circ$=xvolt, Keq =K1
2A+2B2C+2D;E$/circ$=yvolt,Keq=K2
--- PAGE 13 ---
400
PROBLEMSINCHEMISIRY
then:
(a) x=y,K=K2
(b) x = 2y,K= 2K2
(c) x= y,K$2$ =K2
(d)x$2$=y,K$2$=K2
\end{problembox}

\begin{problembox}
\textbf{Problem 36} \hfill \textbf{[Narendra Avasthi Q56]}
\label{q:ext-narendra_avasthi_all_questions-36}

Which graph correctly correlates Ece as a function of concentrations for the cell
Zn(s)+ 2Ag+(aq) -$/to$Zn+(aq)+ 2Ag(s),Ecel =1.56 V
Y-axis : Eceu, X-axis : log10 
[Zn$2$]
[Ag+]
1.56V
.56V
(a)
(b)
0
1.0
1.0
0
1.0
1.56V
(c)
(d)
1.56V
-1.0
0
1.0
-1.0
0
1.0.
\end{problembox}

\begin{problembox}
\textbf{Problem 37} \hfill \textbf{[Narendra Avasthi Q57]}
\label{q:ext-narendra_avasthi_all_questions-37}

The Nernst equation E =E$/circ$ -RT/nF InQ indicates that the Q will be equal to equilibrium
constant K. when:
(a)E=E$/circ$
(b) RT/nF =1
(c) E =zero
(d) E$/circ$=1
\end{problembox}

\begin{problembox}
\textbf{Problem 38} \hfill \textbf{[Narendra Avasthi Q58]}
\label{q:ext-narendra_avasthi_all_questions-38}

The cell reaction 2Ag+(aq)+ H2(g)-$/to$ 2H+(aq)+ 2Ag (s), is best represented by :
(a)Ag(s)|Ag+(aq)|H+(aq)|H2(g)|Pt(s)(b) Pt(s)|H2(g)|H+(aq)| Ag+(aq)|Ag(s)
(c) Ag(s)|Ag+(aq)|| H2(g)| H+(aq)| Pt(s) (d) Ag+(aq)|Ag(s)|| H2(g)| H+(aq)
 + +((+( 
by: 
(a) Cu(s)|Cu+2(aq) II Hg2Cl2(s)| Hg (l)
(b) Cu(s)|Cu+2(aq).I| Hg(l)|HgCl2(s)
(c) Cu (s)| Cu+2(aq)1ICl- (aq)| Hg2Cl2(s)| Hg(l)|Pt(s)
(d) Hg2Cl2(s)|Cl"(aq) |I Cu+2(aq)| Cu(s)
\end{problembox}

\begin{problembox}
\textbf{Problem 39} \hfill \textbf{[Narendra Avasthi Q60]}
\label{q:ext-narendra_avasthi_all_questions-39}

The cellreaction Cr2O2- (aq)+14H+(aq)+ 6Fe$2$+(aq) -$/to$ 6Fe3+(aq)+ 2Cr3+(aq)+7H2O(l),
is best represented by :
(a)Pt(s)|Fe+2(aq),Fe3+(aq) |1Cr2O$2$(aq),Cr3+(aq) |Pt(s)
(b) Pt(s)| Cr2O$2$-(aq),Cr+3(aq) I1 Fe3+(aq),Fe+2(aq)| Pt(s)
(c) Fe$2$+(aq) |Fe+(aq) ICr2O$2$(aq)| Cr3+(aq)
(d) Cr2O$2$(aq)|Cr3+(aq)||Fe3+(aq)|Fe2+(aq)
\end{problembox}

\begin{problembox}
\textbf{Problem 40} \hfill \textbf{[Narendra Avasthi Q61]}
\label{q:ext-narendra_avasthi_all_questions-40}

Select the correct cell reaction of the cell Ag(s)|Ag*(aq) li Cu$2$t(aq)|Cu(s):
(a) 2Ag(s)+Cu(s)---$/to$Cu+2(aq).+ 2Ag+(aq)
(b) Cu(s)+ 2Ag+(aq)---$/to$Cu$2$+(aq) + 2Ag(s)
--- PAGE 14 ---
401
(c) 2Ag (s)+Cu2+(aq) --$/to$ Cu(s)+ 2Ag+(aq)
(d) Cu$2$+(aq)+ 2Ag+(aq)----$/to$ 2Ag(s)+Cu(s)
62.Select the correct cell reaction of the cell Pt(s)|Cl2(g)|Cl-(aq)||Ag+(aq)|Ag(s):
(a) Cl2(g)+Ag+(aq) --$/to$ Ag (s)+ 2C1-(aq)
(b) Cl2(g)+Ag(s)---$/to$2Cl-(aq)+Ag+(aq)
(c)2Cl-(aq)+2Ag+(aq)---$/to$2Ag(s)+Cl2(g)
(d) AgCl(s)----$/to$Ag+(aq)+Cl-(aq)
\end{problembox}

\begin{problembox}
\textbf{Problem 41} \hfill \textbf{[Narendra Avasthi Q63]}
\label{q:ext-narendra_avasthi_all_questions-41}

Standard electrode potential of SHE at 298K is:
(a) 0.05V
(b) 0.10 V
(c) 0.50 V
(d)0.00v
\end{problembox}

\begin{problembox}
\textbf{Problem 42} \hfill \textbf{[Narendra Avasthi Q64]}
\label{q:ext-narendra_avasthi_all_questions-42}

The e.m.f. of the following galvanic cells :
(a) Zn|Zn$2$ (1 M)IICu2+(1M)|Cu
(b) Zn|Zn2+ (0.1 M)IICu2+ (1 M)|Cu
(c) Zn|Zn$2$(1 M)||Cu$2$+(0.1 M)|Cu
(d) Zn|Zn2+(0.1 M) |ICu2+ (0.1 M)|Cu
are represented by E,E2,Eand Erespectively.Which of the following statement is true?
(a)E>E2>E3>E4
(b)E3>E2>E>E4
(c)E>E=E>E2
(d)E2>E=E>E3
65.Based on the cell notation for a spontaneous reaction,at the anode:
Ag(s)|AgCl(s)|Cl-(aq)||Br-(aq)|Br2(l)|C(s)
(a)AgCl becomes reduced
(b)Ag becomes oxidized
(c)Br-becomes oxidized
(d) Br2 becomes reduced
66.Given the listed standard electrode potentials,what is E$/circ$for the cell:
4BiO+(aq) + 3N2H(aq) $/to$ 4Bi(s) + 3N2(g)+ 4H2O(I) + 7H+(aq)
N2(g)+ 5H+(aq)+4e$2$----$/to$NH(aq),E$/circ$=-0.23 V
BiO+(aq)+2H+(aq)+3e$2$---$/to$Bi(s)+H2O(l), E$/circ$=+0.32V
(a)+0.55
(b)+0.34
(c) +1.88
(d)+0.09
67.What is the standard electrodepotential for thereduction of HClO?
HClO(aq)+H+(aq)+ 2e$2$---$/to$Cl-(aq)+H2O(l)
Given:
Cr$2$(aq)$/to$Cr$2$+(aq)+e", E$/circ$=0.41V
HClO(aq)+ H+(aq)+ 2Cr$2$(aq)---$/to$ 2Cr$2$(aq) +CI-(aq) +H2O(l),E$/circ$=1.80
(a) 1.39
(b) 1.54
(c)1.22
(d) 0.90
68.The E$/circ$for the following cell is +0.34 V. In(s)|In(OH)3(aq)l|SbO2(aq)|Sb(s).
Using E$/circ$=-1.oV for the In(OH)31In couple,calculate E$/circ$for the SbO|Sb half-reaction:
(a)-1.34
(b) + 0.66
(c) +0.82
(d) -0.66
69.From the following half-cell reactions and their potentials, what is the smallest possible
standard e.m.f. for spontaneous reactions?
PO(aq)+ 2H2O(l) + 2e$2$---$/to$HPO2 + 3OH-(aq);E$/circ$=-1.05 V
PbO2(s) + H2O(l) + 2e----$/to$PbO(s) + 2OH-(aq);E$/circ$ =+ 0.28 V
--- PAGE 15 ---
402
IO(aq)+2H2O(l)+ 4e-$/to$ IO-(aq) + 4OH"(aq);E$/circ$=+0.56 V
(a) +0.00
(b) +0.74
(c)+0.56
(d) +0.28
(a) HPO2
(b) PO$2$
(c) 10-
(d) 103
\end{problembox}

\begin{problembox}
\textbf{Problem 43} \hfill \textbf{[Narendra Avasthi Q71]}
\label{q:ext-narendra_avasthi_all_questions-43}

Which substance is the best oxidizing agent in Q. no. 47?
(a) 103
(b) 1O-
(c)PbO
(d) PO3
\end{problembox}

\begin{problembox}
\textbf{Problem 44} \hfill \textbf{[Narendra Avasthi Q72]}
\label{q:ext-narendra_avasthi_all_questions-44}

Consider the following half-cell reactions and associated standard half-cell potentials and
determine the maximum voltage that can be obtained by combination resulting in
spontaneous processes:
AuBr (aq)+3e$/to$ Au(s)+ 4Br(aq); .E$/circ$=-0.86 V
Eu3+(aq)+e-$/to$Eu$2$+(aq);E$/circ$=-0.43V
Sn$2$+(aq)+2e$2$---$/to$Sn(s);
E$/circ$=-0.14V
IO-(aq)+H2O(l)+2e---$/to$I(aq)+2OH(aq);E$/circ$=+0.49 V
(a)+0.72
(b)+1.54
(c) +1.00
(d) +1.35
\end{problembox}

\begin{problembox}
\textbf{Problem 45} \hfill \textbf{[Narendra Avasthi Q73]}
\label{q:ext-narendra_avasthi_all_questions-45}

The position of some metals in the electrochemical series in decreasing electropositive
character is Mg > Al > Zn > Cu >Ag. What will happened if copper spoon is used to stirred a
solution of aluminium nitrate?
(a) The spoon gets coated with aluminium.
(b) An alloy of aluminium and copper is formed.
(c)No reaction occurs
(d) The solution starts turning blue
74.Zn can displace:
(a) Mg from its aqueous solution
(b) Cu from its aqueous solution
(c) Na from its aqueous solution
(d) Al from its aqueous solution
\end{problembox}

\begin{problembox}
\textbf{Problem 46} \hfill \textbf{[Narendra Avasthi Q75]}
\label{q:ext-narendra_avasthi_all_questions-46}

Based on the following information arrange four metals A, B, C and D in order of decreasing
ability to act as reducing agents :
(1) Only A,B and C react with 1 MHCl to give H2(g)
(I) When C is added to solutions of the other metal ions,metalic B and D are formed
(II) Metal C does not reduce An+.
(a) C >A>B>D
(b)C>A>D>B
(c)A >C >D>B
(d) A >C >B>D
76.When an aqueous solution of CuSO4 is stirred with a silver spoon then:
(a) Cu+ will be formed
(b) Ag+ will be formed
(c) Cu2+ will be deposited
(d) None of these
77.. Based on the following information arrange four metals, A, B, C and D in order of increasing
ability to act as reducing agents :
(1) Only C react with 1 M HCl to give H2(g)
(Il)When A is added to solution of the other metal ions, metallic D is formed but not
BorC
(a)'D<A<C<B
(b)A<D<C<B
(c)B<D<A<C
(d)D<A<B<C
\end{problembox}

\begin{problembox}
\textbf{Problem 47} \hfill \textbf{[Narendra Avasthi Q78]}
\label{q:ext-narendra_avasthi_all_questions-47}

In the reaction :
4Fe(s)+ 302(g)------$/to$ 4Fe3+(aq) + 60$2$-(aq)
--- PAGE 16 ---
ELECTROCHEMISTRY
403
which of the following statement is incorrect?
(a)A redox reaction
(b) Fe is reducing agent
(c) O2 is an oxidizing agent
(d) Fe is reduced to Fe3t
79.Which of the following is displaced by Fe?
(a)Ag
(b) Zn
(c) Na
(d)All of these
\end{problembox}

\begin{problembox}
\textbf{Problem 48} \hfill \textbf{[Narendra Avasthi Q80]}
\label{q:ext-narendra_avasthi_all_questions-48}

The standard potential at 25oC for the following half reactions is given :
9-=u-++u
Mg$2$++2e$/to$Mg;E$/circ$=-2.37 V
When Zinc dust is added to the solution of MgCl2.
(a) ZnCl2is formed
(b) Mg is precipitated
(c) Zn dissolved in the solution
(d) No reaction takes place
81.The element which can displace three other halogens from their compound is:
(a) F
(b) C1
(c) Br
(d) 1
\end{problembox}

\begin{problembox}
\textbf{Problem 49} \hfill \textbf{[Narendra Avasthi Q82]}
\label{q:ext-narendra_avasthi_all_questions-49}

Using the standard half-cell potential listed, calculate the equilibrium constant for the
reaction:
Co(s)+ 2H+(aq)-$/to$Co$2$+(aq)+H2(g) at 298 K
Co$2$+(aq) +2e----$/to$Co(s) E$/circ$=-0.277 V
(a) 2.3x 109
(b) 4.8x 104
(c) 4.8x 10?
(d) 4.8 x 1011
H2(g)+2AgCI(s)--->2Ag(s) + 2HCl(aq)
(a) 2.8x 107
(b) 5.2x108
(c) 5.2$/times$ 106
(d) 5.2x103
\end{problembox}

\begin{problembox}
\textbf{Problem 50} \hfill \textbf{[Narendra Avasthi Q84]}
\label{q:ext-narendra_avasthi_all_questions-50}

Electrode potential of the half cell Pt(s)|Hg (l)| Hg2Cl2(s)|Cl"(aq) can be increased by :
(a)Increasing[Cl]
(b) Decreasing [Cl-]
(c) Increasing Hg2Cl2(s)
(d) Decreasing Hg(l)
\end{problembox}

\begin{problembox}
\textbf{Problem 51} \hfill \textbf{[Narendra Avasthi Q85]}
\label{q:ext-narendra_avasthi_all_questions-51}

The equilibrium constant for the following general reaction is 1030. Calculate E$/circ$for the cell at
298 K.
2X2(s)+ 3y2+(aq)---$/to$ 2x2+(aq)+ 3Y(s)
(a)+0.105 V
(b)+0.2955V
(c)0.0985V
(d) -0.2955 V
86.A solution containing H* and D+ ions is in equilibrium with a mixture of H2 and D2 gases at
25$/circ$$/circ/text(C)$.If the partial pressures of both gases are 1.0 atm,find the ratio of [D+]/[H*]:
=-0.003 V)
(a) 1.23
(b) 1.12
(c) 0.11
(d) 1.0
\end{problembox}

\begin{problembox}
\textbf{Problem 52} \hfill \textbf{[Narendra Avasthi Q87]}
\label{q:ext-narendra_avasthi_all_questions-52}

The E$/circ$ at 25C$/circ$ for the following reaction is 0.55 V. Calculate the $/Delta$G$/circ$ in kJ :
4BiO*(aq)+3N2H---$/to$4Bi(s)+3N2(g)+4 H2O(l)+7H+
(a)-637
(b) -424
(c)-106
(d)-318.5
88.Use the following E$/circ$for the electrode potentials, calculate $/Delta$G$/circ$in kJ for the indicated
reaction:
(bD)+H8+(b)OuW+(bD)+S(1)O$2$H+(bD)+uW+(b)+S
MnO4(aq)+ 8H+(aq) + 5e$2$----$/to$Mn$2$+(aq)+ 4H2O(l);E$/circ$=+1.51 V
--- PAGE 17 ---
404
Ce4+(aq) +e  $/to$Ce3+(aq)
E$/circ$=+1.61V
(a)-9.65
(b)-24.3
(c)-48.25
(d) -35.2
89.Consider an electrochemical cell in which the following reaction occurs and predict which
changes will decrcase the cell voltage :
Fe$2$(aq)+Ag+(aq)-$/to$Ag(s)+Fe3+(aq)
(1)decrease the[Ag*]
(I) increase in [Fe]
(l1l)increase the amount of Ag
(a) 1
(b)II and III
(c)II
(d)Iand II
90.Consider the following equation for an electrochemical cellreaction.Which of the following
changes in condition will increase the cell voltage?
H2(g)+PbCl2(s)---$/to$Pb(s)+2HCl(aq)
(1)dissolve concentrated HClOqin the cell solution
(Il) increase the pressure of H2(g)
(II) increase the amount of Pb(s)
(a) III
(b) I and II
(c) II and III
(d)11
\end{problembox}

\begin{problembox}
\textbf{Problem 53} \hfill \textbf{[Narendra Avasthi Q91]}
\label{q:ext-narendra_avasthi_all_questions-53}

The standared clectrode potential for the following reaction is +1.33 V. What is the potential
at pH=2.0?
Cr2O2 (aq, 1M) + 14H+(aq)+ 6e$2$ $/to$ 2Cr3+(aq,1M) + 7H2O(l)
(a)+1.820V
(b) +1.990 V
(c) +1.608 V
(d) +1.0542 V
\end{problembox}

\begin{problembox}
\textbf{Problem 54} \hfill \textbf{[Narendra Avasthi Q92]}
\label{q:ext-narendra_avasthi_all_questions-54}

The standard electrode potential for the following reaction is -0.57 V. What is potential at
..pH=12.0?
TeO2 (aq,1M) + 3 H2O(l)+ 4e----$/to$$/to$Te(s)+ 6OH- (aq)
(a)-0.17V
(b)+0.21V
(c) --0.39 V
(d) +1.95 V
\end{problembox}

\begin{problembox}
\textbf{Problem 55} \hfill \textbf{[Narendra Avasthi Q93]}
\label{q:ext-narendra_avasthi_all_questions-55}

Co|Co$2$ (C2)lICo$2$t(C,)Co; for this cell, $/Delta$G is negative if :
(a) C2 >C
(b) C,>C2
(c)C=C2
(d)unpredictable
\end{problembox}

\begin{problembox}
\textbf{Problem 56} \hfill \textbf{[Narendra Avasthi Q94]}
\label{q:ext-narendra_avasthi_all_questions-56}

What will be the emf for the given cell?
Pt|H(g,P)|H+(aq)||H(g,P)|Pt
(a)
(c)
(d) None of these
F
P2
2F " P2
F"R
95.For the electrochemical cell Pt(s)|H2(g)|H*(1M)ICu2+(1M)|Cu(s), which one of the
1 atm
following statements is true?
(a)H liberated at anode and Cu is deposite at cathode.
(b)Hliberated at cathode and Cu is deposite at anode.
(c) Oxidation occurs at cathode.
(d) Reduction occurs at anode.
\end{problembox}

\begin{problembox}
\textbf{Problem 57} \hfill \textbf{[Narendra Avasthi Q96]}
\label{q:ext-narendra_avasthi_all_questions-57}

In a concentration cell the same reactants are present in both the anode and the cathode
compartments,but at different concentrations.Calculate the emf of a cell containing 0.040 M.
Crt in one compartment and 1.0M Cr3 in the other if Cr electrodes are used in both.
(a)0.028V
(b)0.249V
(c)0.083V
(d) 0.125V
\end{problembox}

\begin{problembox}
\textbf{Problem 58} \hfill \textbf{[Narendra Avasthi Q97]}
\label{q:ext-narendra_avasthi_all_questions-58}

A 1.0 M solution of Cd? is added to excess iron and the system is allowed to reach
equilibrium:What is the concentration of Cd2+?
.Cd$2$(aq)+Fe(s)----$/to$Cd(s)+Fe2+(aq);E$/circ$=0.037
--- PAGE 18 ---
405
T
(a) 0.195
(b) 0.097
(c)0.053
(d)0.145
98.The measured voltage for the reaction with the indicated concentrations is 1.50vCalculate E$/circ$
Cr(s) + 3Ag+ (aq, 0.10 M) $/to$ Ag(s) + Cr3*(aq, 0.30 M)
(a) 1.35
(b) 1.40
(c) 1.65
(d) 1.55
\end{problembox}

\begin{problembox}
\textbf{Problem 59} \hfill \textbf{[Narendra Avasthi Q99]}
\label{q:ext-narendra_avasthi_all_questions-59}

Calculate the standard voltage that can be obtained from an ethane oxygen fuel cellat 25$/circ$$/circ/text(C)$.
C2H(g)+7/2O2(g)---$/to$>2CO2(g)+3H2O(l);$/Delta$G$/circ$=-1467kJ
(a) +0.91
(b) +0.54
(c)+0.72
(d) +1.08
$/cdot$to be 0.7714 V If E.
\underline{\hspace{1cm}}=0.535 V,find the pH of H+H2 half-cell.
(a) 1
(b) 3
(c) 5
(d) 7
101.What is the SRPfor S$2$-|CuS|Cuhalf-cell?
Given:
Ksp(CuS) =10-35 and
(a) 1.034V
(b) 1.0 V
(c)-0.694V
(d) 0.694V
102.Given the following standard electrode potentials, the Ks, for PbBr2 is :
PbBr2(s)+2e$/to$Pb(s)+2Br/textasciitilde()(aq);E$/circ$=-0.248V
Pb2+(aq) + 2e -$/to$ Pb(s);
E$/circ$=--0.126V
(a) 7.4x 10-5
(b) 4.9 $/times$ 10-14
(c) 5.2$/times$10-6
(d) 2.3 $/times$ 10-13
103.The standard free energy change for the following reaction is -210 kJ.What is the standard
cell potential?
2H2O2(aq)-$/to$2H2O(1)+O2(g)
(a)+0.752
(b) +1.09
(c)+0.420
(d) +0.640
\end{problembox}

\begin{problembox}
\textbf{Problem 60} \hfill \textbf{[Narendra Avasthi Q104]}
\label{q:ext-narendra_avasthi_all_questions-60}

At equilibrium :
(a) En =0, $/Delta$G= 0
0= 9=
(c)both are correct
(d) none is correct
105.The E$/circ$at 25$/circ$C for the following reaction at the indicated concentrations is1.50 V Calculate
the $/Delta$G in kJat 25$/circ$C:
Cr(s)+ 3Ag+(aq,0.1 M)---$/to$Ag(s) +Cr$2$+(aq,0.1 M)
(a) -140.94
(b)-295
(c)-212
(d)-422.83
106.IfE$/circ$
is 1.69 V and E
Au*/Au
Au3+/Au
(a)0:19V
(b)2.945 V
(c)1.255V
(d)None of these
107.Consider the following standard electrode potentials and calculate the equilibrium constanf at
25$/circ$Cfor the indicated disproportionation reaction:
3Mn$2$(aq)---$/to$Mn(s)+2Mn(aq)
Mn$2$+(aq)+e---$/to$Mn$2$+(aq);E$/circ$=1.51V
Mn$2$+(aq)+ 2e"$/to$Mn(s);
E$/circ$=-1.185V
(a) 1.2 $/times$ 10-43
(b) 2.4 $/times$ 10-73
(c) 6.3$/times$ 10-92
(d) 1.5 $/times$ 10-62
--- PAGE 19 ---
406
PROBLEMSINCHEMISTRY
108.A galvanic cell is composed of two hydrogen electrodes,one of which is a standard one.In
which of the following solutions should the other electrode be immersed to get maximum
e.m.f. :
(a) 0.1 M HCl
(b)0.1MH2SO4(c)0.1M NH4OH
(d)0.01MHCOOH
(a) C=C2
(b) C,>C2
(c) C2 >C
(d) Both (a)and (c)
110.By how much is the oxidizing power ofCr2O2|Cr3+couple decreased if the H+ concentration
is decreased from1 M to10-3M at 25$/circ$C?
(a) 0.001V
(b) 0.207V
(c) 0.441 V
(d) 0.414 V
111.The temperature coefficient of a cell whose operation is based on the reaction
Pb(s)+HgCl2(aq)-$/to$ PbCl2(aq) +Hg(l) is :
E
=1.5$/times$10$/to$vK-1 at 298 K
dT
The change in entropy (in J/K mol) during the operation is :
(a)8627
(b) 57.9
(c)28.95
(d) 14.475
112.Thermodynamic efficiency of a cell is given by :
(a)H
(b) nFE
(c)=nEF
(d) nFEo
AG
$/Delta$G
$/Delta$H
\end{problembox}

\begin{problembox}
\textbf{Problem 61} \hfill \textbf{[Narendra Avasthi Q113]}
\label{q:ext-narendra_avasthi_all_questions-61}

Calculate the value of equilibrium constant (K )for the reaction :
Zn$2$+(aq)+ 4OH"(aq) Zn(OH)$2$(aq)
Given :
Zn$2$(aq)+ 2e----$/to$Zn(s);E$/circ$=-0.76 V
Zn(OH)$2$(aq)+2e---$/to$Zn(s)+4OH-(aq);E$/circ$=-1.36V
=0.06
F
(a)1010
(b)2$/times$1010
(c)1020
(d)none of these
114.Which of the following statement is false for.fuel cells?
(a)They aremoreefficient
(b)They are free frompollution
(c)Theyrun till reactants are active
(d) Fuel burned with O2
115.When a lead storage battery is charged it acts as:
(a)a fuel cell
(b)an electrolytic cell
(c) a galvanic cell
(d)a concentration cell
116.The metal that forms a self-protecting film of oxide to prevent corrosion is :
(a)Na
(b) Al
(c) Cu
(d) Au
117.Rusting of iron is catalyzed by.which of the following?
(a)Fe
(b) Zn
()O2
(d)H+
118.Which of the following is a highly corrosive salt ?
(a) Hg2Cl2
(b)HgCl2
(c)FeCl2
(d)PbCl2
\end{problembox}

\begin{problembox}
\textbf{Problem 62} \hfill \textbf{[Narendra Avasthi Q119]}
\label{q:ext-narendra_avasthi_all_questions-62}

The Zn acts as sacrificial or cathodic protection to prevent rusting of iron because :
(a) Eop of Zn <Eop of Fe
(b) Eop of Zn >Eop of Fe
(c) Eop of Zn =Eop of Fe
(d)Zn is cheaper than iron
--- PAGE 20 ---
407
120.In electrochemical corrosion of metals,the metal undergoing corrosion:
(a) acts as anode
(b) acts as cathode
(c) undergoes reduction
(d)None
\end{problembox}

\begin{problembox}
\textbf{Problem 63} \hfill \textbf{[Narendra Avasthi Q121]}
\label{q:ext-narendra_avasthi_all_questions-63}

When an acid cell is charged, then :
(a) Voltage of cell increases
(b)Resistance of cell increases
(c) Electrolyte of cell dilutes
(d)None of these
122.Electrolytic conduction is due to the movement of:
(a)electrons
(b) ions
(c)atoms
(d) electrons as well as ions
1.23.Molten sodium chloride conducts electricity due to the presence of:
(a)Free electron
(b) Free ions
(c)Free molecules
(d) Atoms of sodium and chlorine
124.Pure water does not conduct electricity because it :
(a)is neutral
(b) is readily decomposed
(c)is almost totally unionized
(d) has a low boiling point
\end{problembox}

\begin{problembox}
\textbf{Problem 64} \hfill \textbf{[Narendra Avasthi Q125]}
\label{q:ext-narendra_avasthi_all_questions-64}

The relation among conductance (G), specific conductance (k) and cell constant (l/A) is :
(a)G=K-
1
(b) G =K
A
(c) GK =
(d) G =KAl
A
A
\end{problembox}

\begin{problembox}
\textbf{Problem 65} \hfill \textbf{[Narendra Avasthi Q126]}
\label{q:ext-narendra_avasthi_all_questions-65}

If x is specific resistance (in S-cm) of the electrolyte solution and y is the molarity of the
solution, then Am (in S cm$2$mol-1) is given by :
(a)
1000x
(b)1000
(c)1000
(d)
xy
y
x
Kx
1000
127.Equivalent conductance can be expressed in terms of specific conductance (k) and
(a)kxN
(b)
K$/times$1000
(c)KxN
(d) K x N x1000
N
1000
\end{problembox}

\begin{problembox}
\textbf{Problem 66} \hfill \textbf{[Narendra Avasthi Q128]}
\label{q:ext-narendra_avasthi_all_questions-66}

Resistance of a decimolar solution between two electrodes 0.02 meter apart and 0.0004 m2 in -
area was found tobe 50 $/Omega$.Specific conductance (k) is:
(a) 0.1 Sm-1
(b) 1S m-1
(c) 10S m-1
(d) 4$/times$10-4 S m-1
129.Resistance of 0.1 M KCl solution in a conductance cell is 300 $/Omega$ and conductivity is
0.013 Scm-1. The value of cell constant is :
(a) 3.9 cm-1
(b) 39 m -1
(c) 3.9 m-1
(d) None of these
130.Ionisation constant of a weak acid (HA) in terms of A and Am is :
CAm
CA2m
(a) Ka =
(b) Ka=
("v -"v)"V
($/alpha$V-"v)
C(Am)2
(c)Ka=
(d)None of these
A(Am-Am)
131.When a concentrated solution of an electrolyte is diluted?
(a) its specific conductance increases
(b)its equivalentconductance decreases
--- PAGE 21 ---
408
(c)its specific conductance decreases and equivalent conductance increases
(d)bothspecificandequivalentconductanceincrease
132.Molar conductivity of a solution of an electrolyte AB is 150 Scm$2$mol-1.If it ionises as
AB3 ---$/to$ A3+ + 3B", its equivalent conductivity will be :
(a) 150 (in Scm$2$eq-1)
(b) 75 (in Scm$2$eq-)
(c) 50 (in Scm$2$eq-1)
(d) 80 (in Scm$2$eq-1)
133.Equivalent conductivity of Fe2(SO4)3 is related to molar conductivity by the expression:
"v= bv (e)
(b) Aeq = Am/3
"v=b$2$v ()
(d) Aeq = Am/6
\end{problembox}

\begin{problembox}
\textbf{Problem 67} \hfill \textbf{[Narendra Avasthi Q134]}
\label{q:ext-narendra_avasthi_all_questions-67}

The limiting equivalent conductivity of NaCl, KCl and KBr are 126.5, 150.0 and 151.5S cm2
eq, respectively. The limiting equivalent ionic conductance for Br--- is 78 S cm$2$eq-1. The
limiting equivalent ionic conductance for Na* ions would be :
(a) 128
(b) 125
(c) 49
(d) 50
135.The specific conductance of a saturated solution of silver bromide is k Scm-. The limiting
ionic conductivity of Ag+ and Br-ions are x and y, respectively. The solubility of silver
bromide in gL-1 is : (molar mass of AgBr = 188)
x(@)
$/times$188
(c)K$/times$1000$/times$188
(d)
x+y
1000
x-y
x+y
x+y
K
188
\end{problembox}

\begin{problembox}
\textbf{Problem 68} \hfill \textbf{[Narendra Avasthi Q136]}
\label{q:ext-narendra_avasthi_all_questions-68}

The resistance of 0.1 N solution of formic acid is 200 $/Omega$ and cell constant is 2.0 cm-1. The
equivalent conductivity (in Scm$2$eq-1) of 0.1 N formic acid is :
(a) 100
(b) 10
(c) 1
(d) none of these
\end{problembox}

\begin{problembox}
\textbf{Problem 69} \hfill \textbf{[Narendra Avasthi Q137]}
\label{q:ext-narendra_avasthi_all_questions-69}

A conductance cell was filled with a 0.02 M KCl solution which has a specific conductance of
2.768$/times$ 10-3 $/Omega(-1)$ cm-. If its resistance is 82.4 $/Omega$ at 25$/circ$C,the cell constant is :
(a)0.2182cm -1
(b) 0.2281 cm -1
(c) 0.2821 cm -1
(d) 0.2381 cm -1
\end{problembox}

\begin{problembox}
\textbf{Problem 70} \hfill \textbf{[Narendra Avasthi Q138]}
\label{q:ext-narendra_avasthi_all_questions-70}

The ionic conductivity of Ba2 and Cl at infinite dilution are 127 and 76 $/Omega(-1)$ cm$2$eq-1
respectively.The equivalent conductivity of BaCl2 at infinity dilution (in $/Omega$/textasciitilde()cm$2$eq-1)
would be :
(a)203
(b) 279
(c) 101.5
(d) 139.5
139.Unit of ionic mobility is :
(a) m V-1 s-1
(b) m$2$v-2 s-1.
(c)m$2$v-1s-1
(d) m-2 V s-1
\end{problembox}

\begin{problembox}
\textbf{Problem 71} \hfill \textbf{[Narendra Avasthi Q140]}
\label{q:ext-narendra_avasthi_all_questions-71}

AAgci can be obtained :
(a)by extraplotation of the graph A and$/sqrt()$C to zero concentration
(b) by known values of A" of AgNO3, HCl and HNO3
(c)both (a) and (b)
(d) None of these
\end{problembox}

\begin{problembox}
\textbf{Problem 72} \hfill \textbf{[Narendra Avasthi Q141]}
\label{q:ext-narendra_avasthi_all_questions-72}

The conductance of a salt solution (AB) measured by two parallel electrodes of area 100 cm2
separated by 10 cm was found to be 0.0001 Q-. If volume enclosed between two electrode
contain 0.1 mole of salt, what is the molar conductivity (Scm$2$mol-l) of salt at same
concentration:
(a) 10
(b) 0.1
(c) 1
(d) none of these
--- PAGE 22 ---
ELECTROCHEMISTR
409
142.Theconductivityofastrongelectrolyte: 
(a) Increases on dilution
(b) Decreases on dilution
(c)Does not change with dilution
(d)Depends upon density of electrolytes
143.The increase in equivalent conductance of a weak electrolyte with dilution is due to:
(a)Increaseindegreeofdissociation and decreaseinionicmobility
(b)Decreasein degreeof dissociation and decrease inionic mobility
(c)Increase in degree of dissociation and increase in ionic mobility
(d) Decrease in degree of dissociation and increase in ionic mobility
144.Strong electrolytes are those which:
(a) Conduct electricity
(b) Dissolve readily in water
(c)Dissociateintoionsathigh dilution
(d) Completely dissociate into ions
145.The electricconduction of a salt solution inwater depends on the:
(a) Size of its molecules
(b)Shape of its molecules
(c)Size of solvent molecules
(d) Extent of its ionization
146.A graph was plotted.between molar conductivity of various electrolytes (NaCl, HCl and
NH4OHI) and $/sqrt()$C (in molL-1). Correct setting:
conductivity
Molar
$/sqrt()$C
(a)I(NaCI), II(HCI),III(NH4OH)
(b) I(HCI), II(NaCI), II(NH4OH)
(C) I(NH4OH), I1(NaCI), III(HCI)
(d) I(NHOH), II(HCI), II(NaCI)
147.Which of the following is arranged inincreasing order of ionicmobility?
(a)I<Br"<Cl-<F"
(b)F-<Cl"<Br<I"
(c)F"<I<Cl<Br"
(d)F-<Cl-<I<Br"
148.HNO3(aq)is titrated with NaOH(aq) conductometrically,graphical representation of the
titration as:
(a)
(b)
8
Vol of NaOH
Vol ofNaOH
(C)
(d)
Vol of NaOH
Vol of NaOH
--- PAGE 23 ---
\end{problembox}

\begin{problembox}
\textbf{Problem 73} \hfill \textbf{[Narendra Avasthi Q149]}
\label{q:ext-narendra_avasthi_all_questions-73}

Which of the following plots willobtained for a conductometric titration of strong acid against
aweakbase?
(a)
(b)
Vofofalkaline
Volofalkaline
uonjos
Volofalkaline
Volof allkaline
solution
\end{problembox}

\begin{problembox}
\textbf{Problem 74} \hfill \textbf{[Narendra Avasthi Q150]}
\label{q:ext-narendra_avasthi_all_questions-74}

Conductometric titration curve of a equimolar mixture ofa HCland HCN with NaOH(aq)is :
(a)
(b)
VolofNaOH
Vol ofNaOH
Conductance
(c)
(d)
Vol of NaOH
Vol of NaOH
--- PAGE 24 ---
ELECTROCHEMISTRY
Level2
1.In the Hall process, aluminum is produced by the electrolysis of molten Al2Oby the following
reactions.Howmany second would it take to produce enough aluminum by theHall process to
make a case of 24 cans of aluminum soft-drink,if each can uses 5.0 g of Al, a current of 9650
amp is employed, and the current efficiency of the cell is 90.0/%:
(a)203.2
(b) 148.14
(c)333
(d) 6.17
2.108 g fairly concentrate solution of AgNOis electrolyzed using 0.1F of electricity.The weight
of resulting solution is:
(a) 94 g
(b) 11.6 g
(c)96.4g
(d)None of these
\end{problembox}

\begin{problembox}
\textbf{Problem 75} \hfill \textbf{[Narendra Avasthi Q3]}
\label{q:ext-narendra_avasthi_all_questions-75}

The electrolysis of acetate solution produces ethane according to reaction :
2CHCOO---$/to$C2Hg(g)+2CO2(g)+ 2e
The current efficiency of the process is 80/%.What volume of gases would be produced at 27$/circ$$/circ/text(C)$
and 740 torr, if the current of 0.5 amp is passed through the solution for 96.45 min?
(a) 6.0 L
(b) 0.60 L
(c) 1.365 L
(d) 0.91 L
4.A layer of chromium metal 0.25 mm thick is to be plated on an auto bumper with a total area
of 0.32 m$2$ from a solution containing CrO2-?What current flow is required for this
electroplating if the bumper is to be plated in 60 s ? The density of chromium metal is
7.20 g/cm3.
(a) 4.9 $/times$ 103A
(b) 1.78 x103A
(c)5.3x104A
(d)10.69$/times$10$/circ$A
\end{problembox}

\begin{problembox}
\textbf{Problem 76} \hfill \textbf{[Narendra Avasthi Q5]}
\label{q:ext-narendra_avasthi_all_questions-76}

100 mL of 0.05 M CuSO4(aq) solution was electrolyzed using inert electrodes by passing
current till the pH of the resulting solution was 2. The solution after electrolysis was
neutralized and then treated with excess KI and formed I2 titrated with 0.04 M Na2S2O3.
Calculate the required volume(inmL) of Na2S2O3:
(a) 112.5 mL
(b) 100 mL
(c)125 mL
(d) None of these
\end{problembox}

\begin{problembox}
\textbf{Problem 77} \hfill \textbf{[Narendra Avasthi Q6]}
\label{q:ext-narendra_avasthi_all_questions-77}

If the equilibrium constant for the reaction H*(aq) + OH-(aq) $/to$ H2O(l) is 103 at certain
temperature then what is the E$/circ$ for the reaction,2H2O(l)+ 2e H2(g)+ 2OH-(aq)
2.303RT
Given:
= 0.066
F
(a) 1.230 V
(b)-0.858V
(c)-0.80V
(d) -0.8274 V
7.Afuel cell develops an electricalpotentialfrom the combustionof butane at 1 bar and 298K
C4H1o(g)+6.5O2(g)--$/to$>4CO2(g)+5H2O(l);$/Delta$,G$/circ$=-2746kJ/mol
What is E$/circ$of a cell?
(a)4.74 V
(b)0.547V
(c) 4.37 V
(d) 1.09 V
\end{problembox}

\begin{problembox}
\textbf{Problem 78} \hfill \textbf{[Narendra Avasthi Q8]}
\label{q:ext-narendra_avasthi_all_questions-78}

The cell Pt|H2(g, 0.1 bar)|H+ (aq), pH = Xl|CI-(aq, 1 M)|Hg2Cl2|Hg|Pt,
has e.m.f. of 0.5755 V at 25$/circ$C.The SOP of caiomel electrode is -0.28V, then pH of solution
will be :
(a) 11
(b) 4.5
(c) 5.5
(d) None of these'
--- PAGE 25 ---
412
\end{problembox}

\begin{problembox}
\textbf{Problem 79} \hfill \textbf{[Narendra Avasthi Q9]}
\label{q:ext-narendra_avasthi_all_questions-79}

For a cell rcaction 2H2(g) + O ,(g)---> 2!12O(l) A,S29s - - 0.32 kJ/K. What is the value of -
$/Delta$fH298 (HO, 1)?
Given:
O2(g)+ 4H+(aq) + 4e---> 2H2O(l);E$/circ$= 1.23 V
(a)-285.07 kJ/mol
(b)-570.14 kJ/mol
(c)285.07 kJ/mol
(d) None of these
10.What is the potential of an electrode which originally contained 0.1 M NO3 and 0.4 M Ht
and which has been treated by 80/% of the cadmiun necessary to reduce allthe NO3 to NO(g)
at 1 bar?
Given:NO3+4H++3e-$/to$NO+2H2O;E$/circ$=0.96V;log 2=0.3
(a) 0.84V
(b) 1.08 V
(c) 1.23 V
(d) 1.36 V
\end{problembox}

\begin{problembox}
\textbf{Problem 80} \hfill \textbf{[Narendra Avasthi Q11]}
\label{q:ext-narendra_avasthi_all_questions-80}

The standard reduction potential of normal calomel electrode and reduction potential of
saturated calomel electrodes are 0.27 and 0.33 voit respectively. What is the concentration of
Cl in saturated solution of KCl?
(a) 0.1 M
(b) 0.01 M
(c)0.001
(d)None
12.Determine the potential of the following cell:
Pt[H(g, 0.1 bar)|H+(aq, 10- M)|MnO4(aq, .1 M),
Mn$2$(aq, 0.01 M),H(aq, 0.01 M)|Pt
Given:E$/circ$
Mno]/ Ma2+ = 1.51 V
(a) 1.54 V
(b) 1.48 V
(c) 1.84 V
(d) none of these
\end{problembox}

\begin{problembox}
\textbf{Problem 81} \hfill \textbf{[Narendra Avasthi Q13]}
\label{q:ext-narendra_avasthi_all_questions-81}

Copper reduces NO into NO and NO2 depending upon concentration of HNO, in solution.
Assuming [Cu2*)= 0.1 M, and Pvo = Pvo2 = 10- bar. At which concentration of HNO3,
thermodynamic tendency for reduction of NO into NO and NO2by copper is same?
E
NO/No = + 0.96 vol,
(a) 101.23 M
(b) 100.56 M
(c) 100.66 M
(d) 100.12 M
\end{problembox}

\begin{problembox}
\textbf{Problem 82} \hfill \textbf{[Narendra Avasthi Q14]}
\label{q:ext-narendra_avasthi_all_questions-82}

For the cell, Pt|Cl2(g, 0.4 bar)|Cl- (aq, 0.1 M)|ICl (aq, 0.01 M)|Cl2(g, 0.2bar)|Pt
The measured potential at 298 K is:
(a) 0.051 V
(b)-0.051 V
(c) 0.102 V
(d)0.0255V
\end{problembox}

\begin{problembox}
\textbf{Problem 83} \hfill \textbf{[Narendra Avasthi Q15]}
\label{q:ext-narendra_avasthi_all_questions-83}

The chlorate ion can disproportionate in basic solution according to reaction,
2ClOCIO+ClO
What is the equilibrium concentration of perchlorate ions from a sclution initially at 0.1 M in
chlorate ions at 298K?
Given:E$/circ$
=0.36 V and E$/circ$
=0.33V at 298 K
(a) 0.019 M
(b) 0.024 M
(c) 0.1 M
(d) 0.19 M
--- PAGE 26 ---
413
16.A cell diagram shown below contains one litre of buffer solution of HA(pK. = 4) and NaA in
both compartment.What is the cell c.m.f.?
H2(g):1bar
H2(e):1bar
iHA (oq, 0.1 M)
HA (aq, 1 M)
and
and
NaA (uq, 1M)
NaA (aq, 1 M)
(a)0.03 V
(b)0.06V
(c) -0.06 v
(d) None of these
\end{problembox}

\begin{problembox}
\textbf{Problem 84} \hfill \textbf{[Narendra Avasthi Q17]}
\label{q:ext-narendra_avasthi_all_questions-84}

Given the cell: Cd(s)|Cd(OH)2(s)|NaOH(aq, 0.01. M)|H2(g,1 bar)| Pt(s)
(a) 0.1
(b) 10-13
(c)10-15
(d) None of these
18.Calculate the e.m.f.(in V) of the cell :
Pt|H2(g)|BOH(aq)|HA(aq)|H(g)|Pt;
0.1 bar 1M
0.1 M
1bar
Given:K. (HA)= 10-7,K,(BOH) = 10-5
(a) 0.39V
(b) 0.36V
(c)0.93V
(d)None of these
\end{problembox}

\begin{problembox}
\textbf{Problem 85} \hfill \textbf{[Narendra Avasthi Q19]}
\label{q:ext-narendra_avasthi_all_questions-85}

Calculate the potential of a half cell having reaction: Ag2S(s)+ 2e 2Ag(s)+ S$2$- (aq)in
a solution buffered at pH=3 and which is also saturated with 0.1 MH2S(aq):
[Given: Ksp(Ag2S) =10-49, Ka $/cdot$Ka = 10-21]
(a) 1.18
(b) 0.39
(c)-0.19 V
(d)none of these
\end{problembox}

\begin{problembox}
\textbf{Problem 86} \hfill \textbf{[Narendra Avasthi Q20]}
\label{q:ext-narendra_avasthi_all_questions-86}

The conductivity of 0.1 N NaOH solution is 0.022 S cm -1. When equal volunne of 0.1 N HCl solution
is added,the conductivity of resultant solution is decreases to 0.0055 S cm"1.The equivalent
conductivity in Scm $2$equivalent/textasciitilde() of NaCl solution is : 
(a)0.0055
(b)0.11
(c)110
(d)None of these
21.In above question after formation of NaCl, further 0.1 N HCl is added,the volume of which is
double to that of the first porrion added, the conductivity increases to 0.018 Scm-. The value
of Aeq (HCl) is [assume no change in conductivity of NaCl (aq)] :
(a) 330S cn$2$eq-1
(b) 305 S cm$2$eq-1
(c) 415 S cm$2$eq-1
(d) 360 s cm2eq-1
\end{problembox}

\begin{problembox}
\textbf{Problem 87} \hfill \textbf{[Narendra Avasthi Q22]}
\label{q:ext-narendra_avasthi_all_questions-87}

Given the following molar conductivities at 25$/circ$C:, HCl, 426Q-1 cm/textasciicircum()mol-1; NaCl,
126Q/textasciitilde()1 cm$2$mol-'; NaC (sodium crotonate), 83Q-'cm$2$mol-1. What is the ionization
constant of crotonic acid? If thc conductivity of a 0.oo1 M crotonic acid solution is
3.83$/times$10-5/textasciitilde()cm-?
(a) 10-5
(b) 1.11 $/times$ 10-5
(c) 1.11 x 104
(d) 0.01
--- PAGE 27 ---
PROBLEMSINCHEMISTRY
\end{problembox}

\begin{problembox}
\textbf{Problem 88} \hfill \textbf{[Narendra Avasthi Q23]}
\label{q:ext-narendra_avasthi_all_questions-88}

Equivalent conductivity of BaCl2, H2SO4 and HCl, are x, X2 and x Scm -eq- at infinite
dilution. If conductivity of saturated BaSO4 solution is x Scm- , then Ksp of BaSO4 is :
500x
106x2
(a)
(b)
(x1 +x2-2xg)
(x1 +x2 -2x)3
2.5$/times$10x2
(c)
0.25x2
(d)
(x1+x2-x3)$2$
(x1+x2-x3)2
\end{problembox}

\begin{problembox}
\textbf{Problem 89} \hfill \textbf{[Narendra Avasthi Q24]}
\label{q:ext-narendra_avasthi_all_questions-89}

The conductivity of 0.001 M Na2SO4 solution is 2.6x 104 Scm- and increases to
7.0 $/times$ 10-4 Scm-1, when the solution is saturated with CaSO4. The molar conductivities of Na*
and Ca2+ are 50 and 120Scm$2$mol-, respectively. Neglect conductivity of used water.
What is the solubility product of CaSO4?
(a) 4$/times$ 10-6
(b) 1.57 $/times$ 10-3
(c) 4x10-4
(d) 2.46 x 10-6
\end{problembox}

\begin{problembox}
\textbf{Problem 90} \hfill \textbf{[Narendra Avasthi Q25]}
\label{q:ext-narendra_avasthi_all_questions-90}

The ionization constant of a weak acid is 1.6x 10-5 and the molar conductivity at infinite
dilution is 380x10-4 S m$2$mol-1. If the cell constant is 0.01 m-1 then conductance of 0.01 M
acid solution is:
(a) 1.52$/times$10-5 s
(b) 1.52 S 
(c)1.52$/times$10-3 s
(d)1.52$/times$10-4 s
26.Three electrolytic cells X,Y,Z containing solution of NaCl, AgNO3 and CuSO respectively
are connected in series combination.During electrolysis 21.6 gm of silver deposite at cathode
in cellY.Which is incorrect statement.
(a)6.35 gm copper deposite at cathode in cell Z
(b) 2.24 litre Cl2 liberated at 1 atm and 273 K at cathode in cell x
(c) 2.24 litre O2 liberated at 1 atm and 273 K at anode in cell Y
(d) 2.24 litre H2 liberated at 1 atm and 273 K at cathode in cell x
moles of H2S2Ogformed is :
(a) 0.25
(b) 0.50
(c) 0.75
(d) 1.00
\end{problembox}

\begin{problembox}
\textbf{Problem 91} \hfill \textbf{[Narendra Avasthi Q28]}
\label{q:ext-narendra_avasthi_all_questions-91}

Zn(s) | Zn(CN)2 (0.5M), CN" (0.01) II Cu(NH)2+ (0.5M), NH(1M)| Cu (s)
Given:
K of Zn(CN)-$2$ =1016
K;of Cu(NH)2=1012,
E$/circ$
0.34V, 
2.303RT
Cu+21Cu
= 0.06
d
The emf of above cell is :
(a) 1.22 V
(b) 1.10 V
(c) 0.98 V
(d) None of these
--- PAGE 28 ---
415
ELEOROCERSTRY
Level 3
PASSOeE
Porous
barrier
Zn(s)
Zn2+NaNOCu2+
Cu(s)
(1M(1M)(1M))
A.Galvanic cell consist of three compartiment as shown in figure. The first compartment
contain ZnSO4(1 M)and ll compartment contain CuSO4(1M).The mid compartment
contain NaNO3(1 M).Each.compartment contain1 L solution:
\end{problembox}

\begin{problembox}
\textbf{Problem 92} \hfill \textbf{[Narendra Avasthi Q1]}
\label{q:ext-narendra_avasthi_all_questions-92}

The concentration of Zn2t in first compartment after passage of 0.1 F charge will be :
(a) 1 M
(b) 1.05 M
(c) 1.025 M
(d) 0.5M
\end{problembox}

\begin{problembox}
\textbf{Problem 93} \hfill \textbf{[Narendra Avasthi Q2]}
\label{q:ext-narendra_avasthi_all_questions-93}

The concentration of NO in mid compartment after passage of 0.1 F of charge will be :
(a)0.95M
(b)0.90 M
(c) 0.975 M
(d).1.05 M
\end{problembox}

\begin{problembox}
\textbf{Problem 94} \hfill \textbf{[Narendra Avasthi Q3]}
\label{q:ext-narendra_avasthi_all_questions-94}

The concentration of SO? ion in Il compartment will be :
(a) 1.05 M
(b) 1.025 M
(c)0.95 M
(d) 0.975 M
.POSSAGE
2
The cell potential (Ece )of a reaction is related as $/Delta$G = -- nF Ece, where $/Delta$G represents max.
usefulelectricalwork
n-no.of moles of electrons exchanged during the reaction
for reversible cell reaction d(AG) =: ($/Delta$,V)dp-.($/Delta$,S). dT
at constant pressure
d(AG)--($/Delta$,S)$/cdot$dT
.At constant pressure.
AG-:AH -T$/cdot$$/Delta$S
..(1)
AG = $/Delta$H + T
(d($/Delta$G)
...(2)
LP
dEll
is known as temperature coefficient of the e.m.f of the cell
dT
--- PAGE 29 ---
416
.1.The temperature coefficient of the e.m.f. of cell,
is given by:
dr
(a)F
b)S
(c)
AS
AS
nF
(d)-nFE
nFT
2.At 300 K, $/Delta$FI for the reaction
Zn(s) +AgCl(s)---$/to$ ZnCl2(aq)+ 2Ag(s) is
-218kJ/mol while the e.m.f.of the cell was1.015V
dE
of the cell is:
(a)4.2 x10-4vk-1
(b)- 3.81 x 10-4vK-1
(c)0.11VK-1
(d) 7.62x 10-4vK-1
3.Colculate$/Delta$S for the given cell reaction in Q.no.2:
(a)-73.53J/Kmol(b)83.53J/Kmol
(c)100J/Kmol
(d) none of these
PASSAGE
3
Molar conductivity(Am)is defined as conducting power of the ions produced by 1 mole of
an electrolyte in a solution.Am =
where k is conductivity (in Scm$2$mol-1)andC is molar
K
C
concentration (in mole/cm$2$)
The molar conductivity of 0.04M solution of MgCl2 is 200 Scm$2$mol at 298 K.A cell with
electrodes that are 2.0 cm$2$ in surface area and 0.50 cm apart is filled with MgCl solution.
1.Conductance of MgCl2 solution is:
(a)8x10-3s
(b)32S
(c)0.032S
(d) None of these'
2.How much current willflow when the potential difference between the two electrodes is 5.0 V?
(a) 156.25 A
(b) 0.16 A
(c)160 A
(d) None of these
PASSAGE
4
In a hydrogen oxygen fuel cell, electricity produced. In this process H2(g)is oxidised at anode
and O2(g)reduced at cathode.
Given : Cathode O2(g) + 2H2O(l)+ 4e/textasciitilde()$/to$ 4OH(aq)
Anode
H2(g)+2OH"(aq)--$/to$2HO(l)+ 2e"
4.48 litre H2 at 1 arm and 273 K oxidised in 9650 sec.
1.The current produced is (in amp) :
(a) 1 amp
(b) 2 amp
(c) 4 amp
(d) 8 amp
\end{problembox}

\begin{problembox}
\textbf{Problem 95} \hfill \textbf{[Narendra Avasthi Q2]}
\label{q:ext-narendra_avasthi_all_questions-95}

The mass of water produced is :
(a) 7.2 gm
(b) 3.6 gm1
(c) 1.8 gm
(d) 0.9 gm
\end{problembox}

\begin{problembox}
\textbf{Problem 96} \hfill \textbf{[Narendra Avasthi Q3]}
\label{q:ext-narendra_avasthi_all_questions-96}

If current produced in fuel cell, use for the deposition of Cu+2 in 1 L, 2MCuSO4 (aq) solution
for 241.25 sec using Pt electrode. The pH of solution after electrolysis is :
(a) 1
(b)2
(c)3
(d) 4
--- PAGE 30 ---
417
ELECTROCHEMISTRY
IPASSAGE
A saturated solution in AgX(Ks,=3x10-12) and AgY(Ksp =10/textasciitilde()12) has conductivity
0.4$/times$10-6Q-cm/textasciitilde()1.
Given : Limiting molar conductivity of Ag+ = 60 Q-cim$2$mol-1
Limiting molar conductivity of X-= 90 Q-cm$2$mol-1
\end{problembox}

\begin{problembox}
\textbf{Problem 97} \hfill \textbf{[Narendra Avasthi Q1]}
\label{q:ext-narendra_avasthi_all_questions-97}

The conductivity of Y- is (in Q- cm-):
(a) 1.45x 10-7
(b) 1.45$/times$10-5
(c) 1.45x 10-9
(d) None of these
2.The limiting molar conductivity of Y- is (in Q- cm$2$mol- ) :
(b)2900
(c) 2.90
(d)None of these
(a)290
ONEORMOREANSWERSIS/ARECORRECT
1.If the e.m.f of a galvanic cell is negative, it implies that :
(a) the cell reaction is spontaneous
.(b)the cell reaction is non-spontaneous
(c)the cell reaction is exothermic
(d) the cell is working in reverse direction
\end{problembox}

\begin{problembox}
\textbf{Problem 98} \hfill \textbf{[Narendra Avasthi Q2]}
\label{q:ext-narendra_avasthi_all_questions-98}

Select right statement(s) about electrolysis :
(b)$/Delta$G is positive for chemical process during electrolysis
(c)Cations and anions are moved toward the anode and cathode respectively
(d) Over voltage is generally associated with evolution of O2 gas
\end{problembox}

\begin{problembox}
\textbf{Problem 99} \hfill \textbf{[Narendra Avasthi Q3]}
\label{q:ext-narendra_avasthi_all_questions-99}

If the half-cell reaction A + e ----$/to$ A- has a large negative reduction potentials, it follows
that:
(a)$/cdot$A is readily reduced
(b) A is readily oxidised
(c) A- is readily reduced
(d)Ais readily oxidised
4.Which of thefollowing statement is correct?
=-0.0V
H+|H2
(a) Cu2+ ions can be reduced by H2(g)
(b) Cu can be oxidized by H+
(c) Sn2+ jons can be reduced by H2
(d) Sn can be oxidized by Cu2+
\end{problembox}

\begin{problembox}
\textbf{Problem 100} \hfill \textbf{[Narendra Avasthi Q5]}
\label{q:ext-narendra_avasthi_all_questions-100}

The oxidation potential of hydrogen half-cell will be negative if :$/cdot$
(a)p(H2)=1atmand[H+]=1M
(b) p(H2) = 1 atm and [H+] = 2 M
(c)p(H2)=0.2atm and [H+] =1 M
(d)p(H2)=0.2 atm and [H+]=0.2 M
6.Which of the following arrangement will produce oxygen at anode during electrolysis?
(a) Dilute H2SO4 with Pt electrodes
(b)Fused NaOH with inert electrodes
--- PAGE 31 ---
418
PROBLEMSINCHEMISTRY
(c) Dilute H2SO4 with Cu electrodes
(d)Concentrate aq.NaCl with Pt electrodes
\end{problembox}

\begin{problembox}
\textbf{Problem 101} \hfill \textbf{[Narendra Avasthi Q7]}
\label{q:ext-narendra_avasthi_all_questions-101}

When an aqueous concentrate solution of lithium chloride is electrolysed using inert
electrodes :
(a) Cl2 is liberated at the anode
(b)Liis deposited at the cathode
(c) as the current flows, pH of the solution around the cathode remains constant
(d) as the current flows,pH of the solution around the cathode increases
\end{problembox}

\begin{problembox}
\textbf{Problem 102} \hfill \textbf{[Narendra Avasthi Q8]}
\label{q:ext-narendra_avasthi_all_questions-102}

Oxygen and hydrogen gas are produced at the anode and cathode during the electrolysisof
fairly concentrate aqueous solution of :
..(a) K2SO4
(b) AgNO3
(c) H2SO4
(d)NaOH
\end{problembox}

\begin{problembox}
\textbf{Problem 103} \hfill \textbf{[Narendra Avasthi Q9]}
\label{q:ext-narendra_avasthi_all_questions-103}

During the purification of copper by. electrolysis :
(a) the.anode used is made of copper ore
(b) pure copper is deposited on the cathode
(c)the impurities such as Ag,Au present in solution as ions
(d) concentration of CuSO4 solution remains constant during dissolution of Cu
\end{problembox}

\begin{problembox}
\textbf{Problem 104} \hfill \textbf{[Narendra Avasthi Q10]}
\label{q:ext-narendra_avasthi_all_questions-104}

When a lead storage battery is discharged :
(a) SO2is evolved
(b) lead sulphate is produced at both electrodes
(c) sulphuric acid is consumed
(d) water is formed
11.Which of the following is characteristic of the cathode in a voltaic cell?
(a)It may gain weight during reaction
(b)Electrons flow to it through the external circuit
(c)It is where oxidation occurs
(d)It received electrons from ions in solution
\end{problembox}

\begin{problembox}
\textbf{Problem 105} \hfill \textbf{[Narendra Avasthi Q12]}
\label{q:ext-narendra_avasthi_all_questions-105}

In an electrochemical process, a salt bridge is used :
(a) tomaintain electrical neutrality in each solution
(b) to complete the external circuit so that current can flow for long time
(c)to mix the solution of anodic and cathodic compartment
(d) to supply.voltage
\end{problembox}

\begin{problembox}
\textbf{Problem 106} \hfill \textbf{[Narendra Avasthi Q13]}
\label{q:ext-narendra_avasthi_all_questions-106}

For a reaction in a galvanic cell the value of -AG$/circ$ at certain temperature is not necessarily
equal to:
(a)nFE
(b)RT In K
(C)T$/cdot$$/Delta$S$/circ$-$/Delta$H
(d)zero
\end{problembox}

\begin{problembox}
\textbf{Problem 107} \hfill \textbf{[Narendra Avasthi Q14]}
\label{q:ext-narendra_avasthi_all_questions-107}

Standard electrode potential of two half-reactions are given below:
Fe$2$ Fe
E$/circ$=-0.44V
Fe3+Fe$2$+
E$/circ$=+0.77V
If Fe?+, Fe?t and Fe are kept together :
(a) the concentration of Fe?t increases
(b) the concentration of Fe?+ decreases
(c)the mass of Fe increases
(d) the concentration of Fe$2$+ decreases
--- PAGE 32 ---
419
ELECTROCHEMISTRY
\end{problembox}

\begin{problembox}
\textbf{Problem 108} \hfill \textbf{[Narendra Avasthi Q15]}
\label{q:ext-narendra_avasthi_all_questions-108}

Which of the following statements are correct regarding to galvanic cell?
(a) A reaction is spontaneous from left to right if Ecen >0
$/cdot$(b) A reaction occurs from right to left if Ecen <0
(c)If the system isat equilibriumnonetreaction occurs
(d) Ece is temperature-independent
\end{problembox}

\begin{problembox}
\textbf{Problem 109} \hfill \textbf{[Narendra Avasthi Q16]}
\label{q:ext-narendra_avasthi_all_questions-109}

Which of the following are concentration cells?
(b) Cd, (Hg)Cd2+|(Hg),Cd
(a)Pt|H2(g)|HC1|H2(g)|Pt
Pi
P2
a(c)a2
nl+nl+uzl(s)uz ()
(d)Ag|AgCl|Cl-(aq)||Br(aq)|AgBr|Ag
C2
C1
C2
\end{problembox}

\begin{problembox}
\textbf{Problem 110} \hfill \textbf{[Narendra Avasthi Q17]}
\label{q:ext-narendra_avasthi_all_questions-110}

In electrolyte concentration cell :
(a) the electrode material and the solution in both half-cells are composed of the same
substances
(b)only the concentrations of solutions of the same substancesis different
(c) Ecel =0
0.0591
(d) the Nernst equation reduces to Ecel
logQ at25$/circ/text(C)$
n
\end{problembox}

\begin{problembox}
\textbf{Problem 111} \hfill \textbf{[Narendra Avasthi Q18]}
\label{q:ext-narendra_avasthi_all_questions-111}

The standard electrode potential of a metal-metal ion (AglAg+) and metal-sparingly soluble
salt anion (AglAgCllCl-)are related as :
RT In Ksp
RT
InKsp
(b)E$/circ$
C1-[AgCI|Ag
=E$/circ$
Ag|Ag
F
RT
Ksp
RT ,[Cl-]
(d)E$/circ$
=E$/circ$
C1|AgCl|Ag
Ag*IAgF
[Cl-]
\end{problembox}

\begin{problembox}
\textbf{Problem 112} \hfill \textbf{[Narendra Avasthi Q19]}
\label{q:ext-narendra_avasthi_all_questions-112}

Which of the following units is correctly matched?
(a) SI units of conductivity is S m-1
.(b) SI units of molar conductivity is S cm$2$mol-1
() SI unit of conductance is S-1
(d)All of these
20.Which of$/cdot$the following statements is/are correct?
(a)The conductance of one cim$2$(or 1 unit3)of a solution is called specific conductance
(b) Specific conductance increases while molar conductivity decreases on progressive
dilution
(c)The limiting equivalent conductivity of weak electrolyte cannot be determine exactly by
extraplotation of theplot ofAeqagainst
(d)The conductance of metals is due to the movement of free electrons
\end{problembox}

\begin{problembox}
\textbf{Problem 113} \hfill \textbf{[Narendra Avasthi Q21]}
\label{q:ext-narendra_avasthi_all_questions-113}

Which is/are correct'statement ?
(a) No corrosion takes place in vacuum
(b) corrosion protecting by electroplating
(c) During rusting Fe2O3$/cdot$ xH2O formed
(d)In presence of electrolyte,corrosion takes place withgreaterrate
22.A dilute solution of KCl was placed between two Pt electrode 10 cm apart across which a
potential difference of 10voltwas applied.Which is/are correct statement(Given:molar
conductivity of K* at infinite dilution is 96.5 Scm$2$ mol-.)
--- PAGE 33 ---
420
Whichis/are correct statement?
(a) Ionic mobility of K+ is 10-3 cm$2$sec volt-1.
(b) The speed of K+ is 10-3 cm sec-1.
(c) Distance travelled by K+ is 5x10$2$ sec is 5 cm.
(d) The potential gradient is 1.0 volt cm -.
(S)()H(D)HOOHWIO()HOHNWIO(H1(S):A
1bar
1bar
pK,(NH4OH)=5; pK. (CHCOOH)= 5;
2.303RT
=0.06
F
Volume of 0.1MNH4OH in anode half cell =100 mL,
Volume of0.1M CHCOOH in cathodehalf cell=100 mL
Which is/are correct statement?
(a) The emf of given cell is 0.48 V.
(b) The emf of given cell is 0.36 V when 50 mL, 0.1 M NaOH added to cathode compartment
(c)The cmf of given cell is 0.36 V when 50 mL 0.1 M HCl added to anode compartment
(d) The emf of given cell is 0.192 V when 100 mL 0.1 M NaOH added to anode compartment
\end{problembox}

\begin{problembox}
\textbf{Problem 114} \hfill \textbf{[Narendra Avasthi Q24]}
\label{q:ext-narendra_avasthi_all_questions-114}

Given : Pt(s)|Cl2(g)|Cl-(C,)I|Cl-(C2)|Cl2(g)| Pt(s)
P) atm
P2atm
Identify in which of following condition working of cell takes place:
(a) C, >C2 and P =P2
(b) P2 >P andC, =C2
(c) C <C2 and P =P2
(d)P2<PandC=C2
25.1000 mL 1M CuSO4(aq) is electrolysed by 9.65 amp current for 100 sec using Pt-electrode.
Whichis/arecorrectstatement?
(a)Blue colour intensity decreases during electrolysis.
(b) Blue colour intensity remains constant if Cu-electrode used.
(c)pH of solution is 8 after electrolysis.
(d) 28 mL of CH4 at 1 atm and 273 K required to its combustion by O2, liberated during
electrolysis.
MATCH THE COLUMN
Column-I and Column-lI containsfourentries each.Entriesof Column-I are tobematchedwith
some entries of Column-Il.One or more than one entries of Column-I may have the matching with
the same entries of Column-II
\end{problembox}

\begin{problembox}
\textbf{Problem 115} \hfill \textbf{[Narendra Avasthi Q1]}
\label{q:ext-narendra_avasthi_all_questions-115}

Column-I
(A)Dilute solution of HCI
(P) O2 evolved at anode
(B)Dilute solution of NaCl
(Q) H2 evolved at cathode
(C)Concentrate solution of NaCl
(R) Cl2 evolved at anode
(D)Fairly concentrate solution of AgNOg
(S) Ag deposition at cathode
--- PAGE 34 ---
421
ELECTROCHENISTRY
Column-!
Column-I
(A)If SOP of substance is exist between$/cdot$
(P) Oxidation of substance is not possible
-1.23 to -0.81 V
(B)If SOP of substance is exist between
,(Q) Oxidation possible only in acidic medium
-0.81 V to -0.40 V
-!
(R) Oxidation possible in any medium
(C) If SOP is less than -1.23 V
(D) If SOP is greater than -0.40 V
.(S) Oxidation easily takes place
Column-I
Column-I (SRP)
3
(P) 0.54
(A)F2+2e=2F-
:(Q) 1.09
(B) Cl2 + 2e/textasciicircum() 2Cl-
(R) 1.36
(C)Br2+2e$2$=Br
(D)I2+2e 2BI
:i(S) 2.87
Column-I
(A) Pt1Fe3, Fe2+
11(P) Metal-metal ion half-cell
(B)Pt|H2|H+
(Q)Gas-gas ion half-cell
(C) Pt|Hg|Hg2
(R) Oxidation-reduction half-cell
(D) Pb|PbSO4|SO2
(S) Metal sparing soluble salt half-cell
Column-I (Property)
Column-II (Unit)
(P) Sm /textasciitilde()1
(A) Conductance
(Q) s-m
(B) Conducivity
$/cdot$(R) Sm$2$mol-1
(C) Molar conductivity
(S) S
(D)Resistivity
Column-I (lon)
Column-II (Molar Conductivity)
H+
'(P)350
(V)
Nat
i(Q) 50
(B)
Li*
(R) 39
(C)
( )
Cs+
:(S) 77
Column-II
Column-1
Galvanic cell
(P) Used in space craft
(V)
(Q)No transformation of electrical energy
(B)
Electrolytic cell
into chemical energy
--- PAGE 35 ---
422
PROBLEMSINCHEMISTRY
(C)Dead battery
(R) Cell reaction is spontaneous
(D) Fuel cell
(S) Cell reaction is non-spontaneous
ASSERTION-REASONTYPEQUESTIONS
Each question contains STATEMENT-1(Assertion)and STATEMENT-2(Reason).
Examine the statements carefully and mark the correct answer according to theinstructions
:- .given below :
(A) If both the statements are TRUE and STATEMENT-2 is the correct explanation of
STATEMENT-1
(B)If both the statements are TRUE but STATEMENT-2is NOT the correct explanation of
STATEMENT-1
(C)If STATEMENT-1is TRUE and STATEMENT-2isFALSE
(D)If STATEMENT-1 is FALSE and STATEMENT-2 is TRUE
\end{problembox}

\begin{problembox}
\textbf{Problem 116} \hfill \textbf{[Narendra Avasthi Q1]}
\label{q:ext-narendra_avasthi_all_questions-116}

STATEMENT-1 : E and n is negative for electrolytic cell.
STATEMENT-2:$/Delta$G$/circ$is +ve for electrolytic cell.
2.STATEMENT-1:When 2 faraday of electricity is passed through 0.1 M H2SO4(aq), 11.2
litre O2evolved at STP
STATEMENT-2:Molecular weight of oxygen is 32.
\end{problembox}

\begin{problembox}
\textbf{Problem 117} \hfill \textbf{[Narendra Avasthi Q3]}
\label{q:ext-narendra_avasthi_all_questions-117}

STATEMENT-1: Copper is dissolved at anode and deposited at cathode when Cu
electrodes are used and electrolyte is 1 M CuSO4(aq) solution.
STATEMENT-2:SOPof Cu isless than SOPof water andSRPof Cuis greater thanSRPof
water.
4.STATEMENT-1:1 coulomb electricity deposits 1 g-equivalent of a substance.
STATEMENT-2:1 faraday is charge on 1 mole of electricity.
5.STATEMENT-1:If SRP of substance is -0.3V, it's reduction is possible at cathode.
STATEMENT-2:Reduction potential of water exist between 0 to-0.8274V at 25$/circ$C.
6.STATEMENT-1:If SRP of substance is-0.SV then reduction of substanceis possible only
inbasic medium.
STATEMENT-2:SRP of water is -0.8274V and at reduction potentialis zero at pH=7
7.STATEMENT-1:The voltage of mercury cell remains constant for longer period of time.
STATEMENT-2:It is because net cellreaction does not involve ions.
8.STATEMENT-1:Lead storage battery is a galvanic cell without salt bridge.
STATEMENT-2:A secondary cell is rechargeable cell.
9.STATEMENT-1:The SRP of three metallic ions A,B,C are -0.3,-0.5,0.8 volt
respectively, so oxidizing power of ions is C>A >B.
STATEMENT-2:Higher the SRP higher the oxidizing power.
10.STATEMENT-1:If SOP of substance is less than -1.23 V and over voltage = 0V,then it's
oxidation in it's aqueous solution is not possible at 298 K.
STATEMENT-2:Standard reduction potential(SRP) of water is +1.23 V.
11.STATEMENT-1:We.cannot add the electrode potentialin order toget electrode potential
.of third electrode if no.of moles of electrons exchanged are not same.
--- PAGE 36 ---
423
ELECTROCHEMISTRY
STATEMENT-2 : Electrode potential is an extensive property.
\end{problembox}

\begin{problembox}
\textbf{Problem 118} \hfill \textbf{[Narendra Avasthi Q12]}
\label{q:ext-narendra_avasthi_all_questions-118}

STATEMENT-1: E and en =O for a chloride ion concentration cell.
RTn
[CI]LHS
nF"
[CI]RHS
dEcell
13.STATEMENT-1:
>O for a cell reaction then $/Delta$S is positive.
IP
E
STATEMENT-2:AS=nFT
\end{problembox}

\begin{problembox}
\textbf{Problem 119} \hfill \textbf{[Narendra Avasthi Q14]}
\label{q:ext-narendra_avasthi_all_questions-119}

STATEMENT-1: Molar conductivity increases with decrease in concentration for weak
electrolytes.
STATEMENT-2: No. of ions increases and no. of ions per unit volume decreases due to
dilution.
and strong electrolytes.
STATEMENT-2:No. of ions per unit volume linearly decreases in both electrolytt.
SUBJECTIVE PROBLEMS
1.How many faradays are required for reductionof 1mol CHsNO2into CHsNH2?
02+2e"-$/to$2OH?
2.What is the equivalent weight ofOin thefollowingreaction,HO+
.2
3.The amount of electricity which releases 2.0g of gold from a gold salt is same as that which
dissolves 0.967g of copper anode during the electrolysisof copper sulphate solution.What is
the oxidation number of gold in the gold state ? (At. wt. of Cu = 63.5; Au =197)
4.When a molten salt was electrolysed for 5 min with 9.65A current, 0.72g of the metal was
deposited.
Calculate the Eq.wt.of metal.
5.During the electrolysis of a concentrated brine solution.Calculate the moles of chlorine gas.
producedby thepassage of 4Felectricity.
6.Calculate the cell potential(in V) if $/Delta$G=-96.5kJ/mol and n =1.
7.If K.for the reaction
Cu$2$+(aq)+Sn$2$+(aq)-$/to$Sn4+(aq)+Cu(s)
at 25$/circ$C is represented as yx109 then find the value of y.
(Given : Ecu2+1Cu
8.If $/Delta$G$/circ$ for the half cell MnO|MnO2 in an acid solution is xF then find the value of x.
(Given : E*
9.If the equilibrium constant for the reaction Cd2+(aq)+ 4NH(aq)= Cd(NH)2(aq) is 10*
then find the value of x.
--- PAGE 37 ---
426
PROBLEMSINCHEMISTRY
Passage-1:
\end{problembox}

\begin{problembox}
\textbf{Problem 120} \hfill \textbf{[Narendra Avasthi Q2]}
\label{q:ext-narendra_avasthi_all_questions-120}

(a)
One or More Answers is/are correct
(a,b,d)
(d)
\end{problembox}

\begin{problembox}
\textbf{Problem 121} \hfill \textbf{[Narendra Avasthi Q10]}
\label{q:ext-narendra_avasthi_all_questions-121}

20.(a,c,d).
21.(a,b,c,d)
\end{problembox}

\begin{problembox}
\textbf{Problem 122} \hfill \textbf{[Narendra Avasthi Q25]}
\label{q:ext-narendra_avasthi_all_questions-122}

(a,b,d)
Match the Column
1.A$/to$P, Q;
B$/to$P, Q;
C$/to$Q,R;
D$/to$P, S
2.A-$/to$Q;
B$/to$R;
C$/to$P;
D$/to$S, R
3.A$/to$S;
B$/to$R;
C$/to$Q;
D$/to$P
4.A$/to$R;
B$/to$Q;
C$/to$P;
D$/to$S
5.A$/to$S;
B$/to$P;
 
D$/to$Q
6.A$/to$P;
C$/to$R;
D$/to$S
7.A$/to$Q,R;
B$/to$S;
C$/to$Q;
D$/to$P,Q,R
Assertion-Reason Type Questions
\end{problembox}

\begin{problembox}
\textbf{Problem 123} \hfill \textbf{[Narendra Avasthi Q15]}
\label{q:ext-narendra_avasthi_all_questions-123}

(C)
Subjective
Problems
\end{problembox}

\begin{problembox}
\textbf{Problem 124} \hfill \textbf{[Narendra Avasthi Q15]}
\label{q:ext-narendra_avasthi_all_questions-124}

7
--- PAGE 38 ---
Hints and Solutions
Level i..
\end{problembox}

\begin{problembox}
\textbf{Problem 125} \hfill \textbf{[Narendra Avasthi Q32]}
\label{q:ext-narendra_avasthi_all_questions-125}

(d) Initial moles ofCr3+ =0.25 $/times$0.2=> 0.05;
6.(c)1F=96500C,
final moles of Cr3+ =0.25 $/times$0.1= 0.025
6.023$/times$1023
moles ofCr3+ reduced,
1C=
=6.24$/times$1018
96500
0.05-0.025=>0.025;
\end{problembox}

\begin{problembox}
\textbf{Problem 126} \hfill \textbf{[Narendra Avasthi Q10]}
\label{q:ext-narendra_avasthi_all_questions-126}

(c)
0.55
0.55$/times$100$/times$60
M
x3=
96500
96500
M = 48.25g/mol
t=75sec
\end{problembox}

\begin{problembox}
\textbf{Problem 127} \hfill \textbf{[Narendra Avasthi Q13]}
\label{q:ext-narendra_avasthi_all_questions-127}

(c)
W
Ixt
0.195
5$/times$60$/times$60xI
34.(a)
W
Ixt
0.838
M
xn=
E
96500'
195
96500
184
xn
96500
I= 21.44
40$/times$60$/times$1.0
\end{problembox}

\begin{problembox}
\textbf{Problem 128} \hfill \textbf{[Narendra Avasthi Q15]}
\label{q:ext-narendra_avasthi_all_questions-128}

(d) When 1 mole of e passed, wt. of substance
96500
deposited
$/to$n=6
38.(d) no.ofeq= no.of F =
1000
1 x 6.023$/times$1023
x8=65.3
-$/to$3.34g=eq.wt
122.5
1.81 $/times$1023
W
1930x0.75
40.(
(a)
W=0.105gm
. atomic wt. of metal =3.34 x3=10.03g
7
96500
8
\end{problembox}

\begin{problembox}
\textbf{Problem 129} \hfill \textbf{[Narendra Avasthi Q41]}
\label{q:ext-narendra_avasthi_all_questions-129}

(a) equivalents of Cl2 produced
16.(c) 1 mole MnO4 lose
x3=10mole
1000$/times$9.65$/times$3600
e;so total charge required=2$/times$10=>20 F
96500
=360
moles of Cl=180
28.(b) Let t sec be used; No.of faradays =
4xt
96500
180$/times$0.0821$/times$300
No. of moles of HO = 4, no. of equivalents
P
1.8
of H0=4$/times$2
=2463L
4xt
\end{problembox}

\begin{problembox}
\textbf{Problem 130} \hfill \textbf{[Narendra Avasthi Q42]}
\label{q:ext-narendra_avasthi_all_questions-130}

(c) No.ofequivalent of H2 produced = Eq.ofO2
96500
=4$/times$2
Ixt
30$/times$193$/times$60
$/to$
t=1.93x105
96500
96500
29.(c) only Au3+,Ag+and Cu2+ will deposit at
=3.6
cathode.
vol. of O2 and H2 produced are
Li will not deposit at cathode because SRP of
=3.6$/times$5.6+3.6x11.2
water is -0.8274 V so after Cu2+;H2 will
= 60.48 litre
evolve at cathode.
\end{problembox}

\begin{problembox}
\textbf{Problem 131} \hfill \textbf{[Narendra Avasthi Q43]}
\label{q:ext-narendra_avasthi_all_questions-131}

(d)1 watt =1 J/sec; 1 kWh =
VxIxt
(in hr)
M
4.6$/times$30$/times$60
30.(d)
1000
M
96500
Total kWh =10/textasciitilde()$/times$110x10$/times$10 = 11
moles ofSn2+ reduced = 0.043
Total cost =11 $/times$2= 22Rs.
Initial moles of Sn2+ = 0.5 $/times$0.6= 0.30;
45.(b) 10 g Cd should present with 90 g Hg.
wt. of Cd required with 4.5 g of Hg at
remaining moles of Sn2+ = 0.30 - 0.043
cathode
0.257
10
[Sn$2$+]=
0.514M
x4.5= 0.5;(Cd$2$+ 2e$/to$Cd)
0.5
90
--- PAGE 39 ---
430
PROBLEMSINCHEMISTRY
146.(b)Ionic molar conductivity of H*is very high
5.(a) Initial m moles ofCu$2$+= 5;
and NHOH is a weak electrolyte.
m-eq. or m-moles of H+ produced
150.(d) Molar conductivity of H+ and OH"are very
=100$/times$10-$2$=1
high as compare to other ions.
m-moies of Cu2+ converted into
Initially conductance of solution sharply
decreases due to consumption of free H*
Cu: 
then increases due to formation of salt
2
(NaCN)and After complete neutralization
m-moles ofCu2+ remaining in solution
further sharply increases due to presence of
=5-0.5=4.5
OH"
2Cu$2$++4I$/to$Cu212+12
Level2
and 12 + 2NaSO3 ---$/to$ 2NaI+ NaSO6
\underline{\hspace{1cm}}Ixnxt
m-moles ofCu2+ remaining = m-moles of
1.(b) No.of equivalent of aluminum,-
W
E
96500
NaS2O3
24$/times$5
9650x0.9xt
4.5=0.04xV
V=112.5mL
27
96500
6.(b)2HO(D2H(aq)+2OH"(aq);$/Delta$G
t =148.14 sec
2H(aq)+2e$2$H(g);
G2
2.(c) AtAnode:
2HO()-$/to$4H*(aq)+O2(g)+ 4e"
2HO()+2eH2(g)+2OH(aq);$/Delta$,G
At Cathode:Ag*(aq)+e---$/to$Ag(s)
$/Delta$,G=$/Delta$G$2$+$/Delta$G2
Eq.ofO2evolved = Eq.of Ag formed= 0.1
-2$/times$F$/times$E$/circ$=-RT1n(10-26)-2xFx0
0.1$/times$32
0.066
Total loss in wt. =
+ 0.1$/times$108
E$/circ$=
log(10-26)
4
2
0.8+10.8=11.6
=-0.858V
wt. of final solution =108 -11.6=> 96.4g
7.(d) To find n we break the cell reaction into two
half cell reduction
96500
Anode:
0.5$/times$0.8$/times$96.5$/times$60
13HO(D---$/to$6.502(g)+26H*(aq.)+26e"
=0.024
96500
Cathode:
moles ofCO2(n =1)produced = 0.024
4CO2(g)+26H*(aq.)+26e$2$$/to$C4Ho(g)
moles of CH (n = 2)produced
+8HO(
0.024
$/Delta$G$/circ$
(-2746)$/times$1000
=0.012
E$/circ$=
= 1.09 V
2
nF
26$/times$96500
Total moles of gases produced = 0.036
\end{problembox}

\begin{problembox}
\textbf{Problem 132} \hfill \textbf{[Narendra Avasthi Q8]}
\label{q:ext-narendra_avasthi_all_questions-132}

(c) For normal calomel electrode Erp = ERp;
nRT
0.036x0.0821x300
740)
0.0591
[H+]$2$
P
Ee =(0.28-0)-
2
PH2
(760)
=0.91litre
0.5755=0.28+0.0591 pH+
0.0591
log (0.1)
4.(d) Total volume of metal layer
2
= 0.25x0.32$/times$10-3
pH=5.5
9.(a)We know
0.08$/times$10/textasciitilde() m$2$
2H2(g)---$/to$4H++4e$2$;E=0.0V
Total wt.of chromium layer
O2(g)+4H*(aq)+4e$2$$/to$2HO();E=1.23V
= 0.08 $/times$7.20x10$/circ$ x10-3= 576g
1$/times$60
2H(g)+O2(g)---$/to$2HO(D; E=1.23V
576
x6=
52
96500
$/Delta$G29g =-nFE$/circ$=-4$/times$96500 x1.23
1=10.69$/times$10A
= - 474.78 kJ
--- PAGE 40 ---
LECIROCREMISTR
$/Delta$H29s =$/Delta$,G29g + T.$/Delta$,S298
=(-474.78) + 298 x(-0.32)
= - 570.14 kJ/mol
$/Delta$H29g = -285.07 kJ/mol
\end{problembox}

\begin{problembox}
\textbf{Problem 133} \hfill \textbf{[Narendra Avasthi Q10]}
\label{q:ext-narendra_avasthi_all_questions-133}

(a) After addition of Cd and it's oxidation into Cd2+
NO(aq)+ 4H*(aq)+3e---$/to$$/to$ NO(g)+ 2HO(D)
0.1-x
0.4 - 4x,
where x=0.08
[NO] remaining=0.02M;[H]remaining = 0.08M
0.0591
E
ONION
3
log
[NO][H+]4
=0.96-
0.0591
log
3
(0.02)(0.08)4
=0.84V
11.(a) E$/circ$ for normal calomel electrode =E$/circ$for saturated calomel electrode
0.33=0.27
0.06
log[C1]$2$;[CI]=0.1M
2
12.(b)Anode:H2(g)-$/to$2H*(aq)+ 2e$/circ$;E$/circ$=0.0V
Cathode:MnO(aq)+8H*(aq)+5e----$/to$Mn$2$+(aq)+ 4H2O(D); E$/circ$=1.51 V
Ece =ERP(RHS)) -ERP(LHS)
0.06
[Mn"+]
log
0.06
5
[MnO][H+
0.06
0.01
1.51
0.06
0.1
801
0.1$/times$(10-2)8
2
log
(10-3)2
=1.48 V
13.(c)3Cu(s)+ 8H*+2NO-
$/circ$ + ON + + <-
Cu(s)+ 4H++ 2NO3---$/to$>Cu2++ 2NO2+ 2H2O
Let concentration ofHNOisX so[H+]=X
X=[ON]pue
E
=E
NONO
ZONICON
or
ONIEON
=E
NO/NO2
0.0591
10-3
0.96
0.0591
3
8o
0.79
1
0.0591
10-9
0.62
0.0591
10-9
log
10
=0.45
log
6
2
x6
log x = 0.657 = 0.66 or x = 10+0.6
V14. (a) Anode :
2Cl$2$(aq)--> Cl(g)+ 2e
Cathode:
Cl2(g)+ 2e$2$---$/to$ 2CI$2$(aq)
V!
Cell reaction:2Cl(aq)+Cl2(g)---$/to$2Cl(aq)+Cl2(g)
FS.
[CI]HS(PaHS
Ecell=
0.06
log
[C1]HS(PCIRHS
(c: Ee = 0)
2
--- PAGE 41 ---
432
PROBLEMSINCHEMISTRY
0.06
(0.01)2
0.4
$/to$0.051V
2
\end{problembox}

\begin{problembox}
\textbf{Problem 134} \hfill \textbf{[Narendra Avasthi Q15]}
\label{q:ext-narendra_avasthi_all_questions-134}

(a)
HO(D)+CIO(aq)-$/to$ClO(aq)+2H+(aq)+ 2e
2H(aq)+ClO(aq)+2e$2$----$/to$2ClO(aq)+HO(D)
2ClO(aq)--$/to$Cl02(aq)+ClO(aq)
RT
Een =0.33-0.36=-0.03;
lnK
2F
0.06
-0.03
logK or K=0.1
2
2Cl0Clo+ClO2
0.1-2x
x
X
x$2$
1
3.16x=0.1-2x
(0.1 --2x)$2$
10'
5.16x = 0.1
x = 0.1/5.16 = 0.0193= 1.9$/times$10-2
0.06
[H*]RHS
\end{problembox}

\begin{problembox}
\textbf{Problem 135} \hfill \textbf{[Narendra Avasthi Q16]}
\label{q:ext-narendra_avasthi_all_questions-135}

(b)
Ecell
1
[H]HS
J0
Ecel = 0.06 [(pH)LHs - (pH)RHs]
(pH)Hs =pKa+ log
[A"]
[HA]
0.1
(pH)RHS = 4
Ecel = 0.06 (5-4)= + 0.06 V
\end{problembox}

\begin{problembox}
\textbf{Problem 136} \hfill \textbf{[Narendra Avasthi Q17]}
\label{q:ext-narendra_avasthi_all_questions-136}

(c)
0.06
[Cd2+]
10
Ecell
2
[H]?
Ecell = 0,
[Cd$2$+][OH-]$2$
Ell = 0.031og >
K2
0.39
log
= 13
K2
0.03
Kp=1013
x(10-14)$2$=10-15
[H*]RHSPHUS
\end{problembox}

\begin{problembox}
\textbf{Problem 137} \hfill \textbf{[Narendra Avasthi Q18]}
\label{q:ext-narendra_avasthi_all_questions-137}

(a)
0.06
Ecell
2
[H2HSPHJRHS
[>>
[OH]=$/sqrt()$K,xC; [OH]=10-3
or
[H*]Hs = 10-11
[H*]RHs =$/sqrt()$KaxC=104
0.06
(104)$2$$/times$0.1
Ecll
801
x1=0.39V
2
(10-11)2
--- PAGE 42 ---
ELECTROC
433
\end{problembox}

\begin{problembox}
\textbf{Problem 138} \hfill \textbf{[Narendra Avasthi Q19]}
\label{q:ext-narendra_avasthi_all_questions-138}

(c)
HS $/to$ 2H+ + s2-
23.(c)A(BaSO4)=2A(BaSO)
[H+]$2$[S2-]
Ka,.Ka2
Aeg(BaSO4)= Aq(Ba$2$+)+ Ag(SO$2$-)
[H,S]
= Aq(BaCl)+Aq(HSO)-Aq(HCI)
=10-21$/times$0.1=(10-3)$2$rs-2]
Aeq(BaSO)=x+x2-x3
[Ss-2] = 10-16
m=2(x+x2-x)
0.06
[s$2$-]
for sparingly soluble salt
Ag* Ag
log
2
Ksp
=K
$/times$1000
10-16
M
=0.80-
0.06
log
or M =
x
2
10-49
2(x+x2-x3)
x1000
0.06
x33=>-0.19V
500x
=0.80-
2
(x+x2-x)
\end{problembox}

\begin{problembox}
\textbf{Problem 139} \hfill \textbf{[Narendra Avasthi Q20]}
\label{q:ext-narendra_avasthi_all_questions-139}

(c) Normality of resultant solution
2.5$/times$10x2
0.1 x V
= 0.05 N
(x+x2-x3)
VxV
\end{problembox}

\begin{problembox}
\textbf{Problem 140} \hfill \textbf{[Narendra Avasthi Q24]}
\label{q:ext-narendra_avasthi_all_questions-140}

(a) Conductivity of Na2SO4 = 2.6 $/times$10-4
0.0055
/textasciicircum()eq =1 $/times$
= 110 S cm$2$eq-1
Am(NaSO4) =
1000$/times$2.6 $/times$10-4
0.05
0.001
\end{problembox}

\begin{problembox}
\textbf{Problem 141} \hfill \textbf{[Narendra Avasthi Q21]}
\label{q:ext-narendra_avasthi_all_questions-141}

(b) Resultant conc. of NaCl =
0.1 xV
0.025N
= 260S cm$2$mol-1
4V
m(SO$2$)=Am(Na2SO)-2m(Na)
andHCl
0.1$/times$2V
= 0.05 N
4V
=260-2$/times$50
= 160 S cm$2$mol-1
Ktotal =KNaCI+KHC1
N(NaCl).Aeq(NaCl)
Conductivity ofCaSOsolution
Ksolution
1000
=7$/times$10--2.6$/times$10-4
N(HCI). Aeq (HCI)
=4.4x10-S cm/textasciitilde()1
 
1000
Am(CaSO)=m(Ca$2$+)+m(SO$2$)
0.025$/times$110
0.05. Aeq(HCI)
=120+160
0.018=
1000
1000
= 280 S cm$2$mol-1
18 -2.75
Solubility C
1000$/times$K
1000 $/times$4.4 $/times$10-4
eg(HCI)
=305 S cm$2$mol-1
0.05
Am
280
\end{problembox}

\begin{problembox}
\textbf{Problem 142} \hfill \textbf{[Narendra Avasthi Q22]}
\label{q:ext-narendra_avasthi_all_questions-142}

(b) The molar conductivity of the dissociated
=1.57 $/times$10-3 M
form of crotonic acid is
Kgp =[Ca$2$+][SO$2$-]otal
Am(HC)= Am(HCl) +Am(NaC)-Am(NaCl)
=(0.00157)(0.00157 + 0.001)
=(426 + 83-126)Q-1cm$2$mol-1
=4.0$/times$10-6M2
=383Q$/Omega$-$2$cm$2$mol-1
25.(b)K. =
C$/alpha$$2$
0.01x$/alpha$$2$
1-$/alpha$
$/to$1.6$/times$10-5=
The molar conductivity of HC,
1-$/alpha$
1.6$/times$10-5
$/times$1000
$/to$0.04
C
0.001
0.01
= 38.3Q- cm$2$mol-1
Am=0.04$/times$380$/times$10-4
The degree of dissociation,
Am
Am(HC)
(38.3Q-1cm$2$mol-1)
=15.2$/times$10-4 S m$2$mol-1
(383Q-cm$2$mol)
0.1
(DH)V
K=AmxC=15.2$/times$10-4$/times$10$2$$/times$10-2
C$/alpha$$2$
(10-)(0.1)$2$
(:1m$2$=1000 litre)
1-0.1
1.11$/times$10-5
k=G.G*and G*=0.01 m-,
--- PAGE 43 ---
434
PROBLEMSINCHEMISTRY
where G*
$/Delta$H=nF
A
(dT)
1.52$/times$10-2
-218$/times$1000=2$/times$96500$/times$300
dE
G=
$/to$1.52S
0.01
27.(a)xH=xo+XHzSOs
-2$/times$96500$/times$1.015
22.4
\underline{\hspace{1cm}}4$/times$8.4
2x
+2$/times$n
-3.81x10-4vK-1
22.4
22.4
2=1.5+ 2n=$/to$n =0.25
3.(a)$/Delta$S=nF
E
= 2$/times$96500$/times$-3.81$/times$10-4
dT
Level 3
=-75.53J/mol-K
Passage--1
Passage-3
1.(c),2.(a),3. (d)
\end{problembox}

\begin{problembox}
\textbf{Problem 143} \hfill \textbf{[Narendra Avasthi Q2]}
\label{q:ext-narendra_avasthi_all_questions-143}

(b) K=Am.C=
m.M
Equivalent ofZn$2$+ produced =0.1 or
1000
moles ofZn2+\underline{\hspace{1cm}} 0.1
0.05
200 $/times$0.04
K=
8$/times$10$2$Scm-
2
1000
+ve charge increases in first compartment so
0.50
due to interaction and maintain electrical
$/to$8$/times$10-3
Zn2+
2
neutrality
move
toward
 I
compartment and NO move towards first
or G=0.032S;V=IR
I
compartment.Solutionis always electrically
G
neutral so charge of1 Zn2+ is neutralized by
I=5x0.032=0.16A
2NO3-
. [Zn2+] in first compartment
SubjectiveProblems
0.05
=1 +
= 1.025 M
14.A(HA)=A(HCI)+A(NaA)-A(NaCI)
2
Concentration
second
= 425+ 100 --125= 400S cm$2$mol-1
Jo
NO3
in
pH=4,[H+]=10-4=$/alpha$C
compartment
=1-0.05=0.95M
200
=0.5;
In third
compartment moles of Cu2+
400
reduced
(Ca)$/cdot$$/alpha$
10-4(0.5)
Ka
=10-;pKa=4
0.05
=0.025
(1-$/alpha$)
(1-$/alpha$)
2
15.E$2$a-/AgCl/Ag=EAg/Ag
0.0591
1ogKsp
Relatively -ve charge increased so SO2 and
1
Nat move toward opposite direction to
0.209=0.80+
0.0591
logKsp
maintain electrical neutrality
1
[SO2]remaining
0.025
>0.975M
Ksp=10-1; Let solubility of AgCl in
2
0.01 M solution is x
Passage-2
10-10 =x(x+ 0.01)
d($/Delta$G)
\end{problembox}

\begin{problembox}
\textbf{Problem 144} \hfill \textbf{[Narendra Avasthi Q1]}
\label{q:ext-narendra_avasthi_all_questions-144}

(b) From equation 1 and 2, $/Delta$S :
x=10-8
dT
d(-nFE)
dE
$/Delta$S
.Moles of AgCl dissolved in
$/Delta$S =
dT
dT
or
dT
nF
10L=10-8$/times$10=10-7
2.(b) From equation 2;
\end{problembox}